% -----------------------------------------------------------
\part{Book of Mormon}
% -----------------------------------------------------------

% -----------------------------------------------------------
\section{1 Nephi}
% -----------------------------------------------------------

	% ----------------------
	\subsection{Chapter 1}
	% ----------------------
	
		\subsubsection{Verse 1}  I, NEPHI, having been born of goodly parents, therefore I was taught somewhat in all the learning of my father; and having seen many afflictions in the course of my days, nevertheless, having been highly favored of the Lord in all my days; yea, having had a great knowledge of the goodness and the mysteries of God, therefore I make a record of my proceedings in my days.
		\subsubsection{Verse 2}  Yea, I make a record in the language of my father, which consists of the learning of the Jews and the language of the Egyptians.
		\subsubsection{Verse 3}  And I know that the record which I make is true; and I make it with mine own hand; and I make it according to my knowledge.
		\subsubsection{Verse 4}  For it came to pass in the commencement of the first year of the reign of Zedekiah, king of Judah, (my father, Lehi, having dwelt at Jerusalem in all his days); and in that same year there came many prophets, prophesying unto the people that they must repent, or the great city Jerusalem must be destroyed.
		\subsubsection{Verse 5}  Wherefore it came to pass that my father, Lehi, as he went forth prayed unto the Lord, yea, even with all his heart, in behalf of his people.
		\subsubsection{Verse 6}  And it came to pass as he prayed unto the Lord, there came a pillar of fire and dwelt upon a rock before him; and he saw and heard much; and because of the things which he saw and heard he did quake and tremble exceedingly.
		\subsubsection{Verse 7}  And it came to pass that he returned to his own house at Jerusalem; and he cast himself upon his bed, being overcome with the Spirit and the things which he had seen.
		\subsubsection{Verse 8}  And being thus overcome with the Spirit, he was carried away in a vision, even that he saw the heavens open, and he thought he saw God sitting upon his throne, surrounded with numberless concourses of angels in the attitude of singing and praising their God.
		\subsubsection{Verse 9}  And it came to pass that he saw One descending out of the midst of heaven, and he beheld that his luster was above that of the sun at noon-day.
		\subsubsection{Verse 10}  And he also saw twelve others following him, and their brightness did exceed that of the stars in the firmament.
		\subsubsection{Verse 11}  And they came down and went forth upon the face of the earth; and the first came and stood before my father, and gave unto him a book, and bade him that he should read.
		\subsubsection{Verse 12}  And it came to pass that as he read, he was filled with the Spirit of the Lord.
		\subsubsection{Verse 13}  And he read, saying: Wo, wo, unto Jerusalem, for I have seen thine abominations! Yea, and many things did my father read concerning Jerusalem—that it should be destroyed, and the inhabitants thereof; many should perish by the sword, and many should be carried away captive into Babylon.
		\subsubsection{Verse 14}  And it came to pass that when my father had read and seen many great and marvelous things, he did exclaim many things unto the Lord; such as: Great and marvelous are thy works, O Lord God Almighty!  Thy throne is high in the heavens, and thy power, and goodness, and mercy are over all the inhabitants of the earth; and, because thou art merciful, thou wilt not suffer those who come unto thee that they shall perish!
		\subsubsection{Verse 15}  And after this manner was the language of my father in the praising of his God; for his soul did rejoice, and his whole heart was filled, because of the things which he had seen, yea, which the Lord had shown unto him.
		\subsubsection{Verse 16}  And now I, Nephi, do not make a full account of the things which my father hath written, for he hath written many things which he saw in visions and in dreams; and he also hath written many things which he prophesied and spake unto his children, of which I shall not make a full account.
		\subsubsection{Verse 17}  But I shall make an account of my proceedings in my days.  Behold, I make an abridgment of the record of my father, upon plates which I have made with mine own hands; wherefore, after I have abridged the record of my father then will I make an account of mine own life.
		\subsubsection{Verse 18}  Therefore, I would that ye should know, that after the Lord had shown so many marvelous things unto my father, Lehi, yea, concerning the destruction of Jerusalem, behold he went forth among the people, and began to prophesy and to declare unto them concerning the things which he had both seen and heard.
		\subsubsection{Verse 19}  And it came to pass that the Jews did mock him because of the things which he testified of them; for he truly testified of their wickedness and their abominations; and he testified that the things which he saw and heard, and also the things which he read in the book, manifested plainly of the coming of a Messiah, and also the redemption of the world.
		\subsubsection{Verse 20}  And when the Jews heard these things they were angry with him; yea, even as with the prophets of old, whom they had cast out, and stoned, and slain; and they also sought his life, that they might take it away.  But behold, I, Nephi, will show unto you that the tender mercies of the Lord are over all those whom he hath chosen, because of their faith, to make them mighty even unto the power of deliverance.


\section{Book of Mormon} % *BOM

\subsection{BOM interesting stories to teach} % *BOMINTERESTING *BOMSTORIES *BOMSTORY

\begin{itemize}

\item \textbf{View Lehi/Nephi/promised land story as a type of our trying to enter into the highest degree of the celestial kingdom (1N)}

	Requires being married (1N 7:1); requires great sacrifice; knowledge from the Lord; death and resurrection (into the water/coming out of it), etc.
	See 1N 8:1, seeds; other verses that mention seed(s).
	Connect to Lehi's vision, the fruit desirable to make one happy? And connect that to 2N 5:27 and 2N 9:43? Lots of other connections to make here, in his dream and the way he describes the fruit and everything.
		So you see Lehi's dream as an allegorical expression of what is happening to his family as they journey to the new land; and then you see both of those as allegorical expressions of our journey to the highest degree of the celestial kingdom.
		Lehi wants his FAMILY to partake of it.
		Water in both the vision and the journey to the new land (the sea).
		Tree in the vision; they take seeds for ``fruit of every kind'' (8:1).
		Lehi doesn't go get them personally after taking the fruit; he beckons and calls to them only. Doesn't approach L/L either; he apparently just beckons and calls but they refuse.
		The role of the SEEDS they carry.
		
\item \textbf{Track the gaining of each person's testimony, and the differences in experiences that each person has as they gain their testimonies (1N)}

	Nephi and Sam, early. Sariah, when her sons return safely. Laman and Lemuel, perhaps earlier, but they say ``we know of a surety'' after Nephi shocks them in 1N 17:52--55, especially 55. What about Zoram? And the members of Ishmael's family? Track each person, very interesting.
	
\item \textbf{In Lehi's dream, how we go through the rod/temptation/fruit cycle over and over, and how in different aspects of our life we are on different places in different paths holding on to different ``rods''}
	
\item \textbf{Compare and contrast the journeys across the water (1N/Ether)}

	Note that Nephi's group goes crazy (18:9) and faces lots of setbacks in the weather (18:13); the boj's group praised the Lord the entire time (Eth 6:9; compare Nephi's attitude in 1N 18:16) but still faced lots of challenges (Eth 6:6--7). The main difference seems to be that the Lord caused the wind to blow the entire time in the boj's favor (Eth 6:5, 8), while with Nephi's group, because of their wickedness I'm guessing, they actually went in reverse, away from their goal (18:13). Note the contrast between 18:15, they're about to die, and Eth 6:10, no monster could break them, and they had light the entire time.
	
	Another good parallel: the boj's people had light the entire time, but for Nephi's group, their source of light was the Liahona, and it stopped working (18:12) until Nephi took charge again (18:21). Since they didn't have the light of the compass, they didn't know where to steer the ship, so they couldn't have even gotten there on their own if they had tried to fight the wind pushing against them.
	
	1N 17:3 kind of describes this situation: when we're righteous, the Lord prepares the way before us.
	
	So a few morals:
		\begin{compactitem}
		\item Nephi had the right attitude, but because of the people around him, he wasn't progressing.
		\item Being righteous isn't a guarantee of calm seas. But it is a guarantee that (1) we'll have the light of the Lord with us, and (2) the wind will always be pushing in the direction of the promised land, whatever happens.
		\item Sometimes, because of the decisions of other people, our prayers are totally ineffectual (1N 18:19).
		\end{compactitem}

\item \textbf{Compare/contrast the hardships and escapes of Limhi's people (Mo 20--22) and those of Alma's people (Mo 23--24)}

	How are they similar? How are they different? Why does the Lord act differently in each case? How do the people respond differently? Note Mo 25:10 as well, which talks about the power of God in delivering them only when mentioning Alma; doesn't mention Limhi's people and the power of God. But then see Mo 25:16.

\item \textbf{The calling and ministry of Amulek}

	He sees a vision of Alma coming and receives him, thinking that Alma will be a blessing to him and his house. But his family ends up abandoning him and he loses everything. Get all the scriptures that talk about him, including where he is talking about himself.

\item \textbf{The conversion of Lamoni as an almost literal rebirth and resurrection (Al 18--19}

	This really starts at the end of 18, where Lamoni falls as if he were dead. He is raised as if from the dead by his wife---what significance does that have? Abish tells everyone the news; he gets up and ``stands'' (common when talking about the resurrection); many beautiful parallels here, including hand-to-hand contact as in the temple and through the veil; see also Al 47:22--24. A beautiful story with a beautiful parallel; this could be powerful when fleshed out. This is almost a literal rebirth for them. Note the word ``ground'' for the grave, parallels like that. The words are very particular. Note what the people say in 19:33, that they have no more desire to do evil. That's what Benjamin's people say in Mo 5:2. Those people didn't fall asleep or go unconscious, but they saw themselves ``in their own carnal state, even less than the dust of the earth.'' There are two cases of death language there---the dust, referring to the grave, and their carnal state, for being carnally minded is death (Rom 8:6, 2N 9:39). So I shouldn't only be looking for the rebirth process in the scriptures, I should also be looking at how that process is a type of the resurrection, and how the language used to describe one is often used to describe the other.
		See also what happens with Lamoni's father, especially in 22:18, 22.
		For more about being dead in sin and raised from it, see Eph 2:1--6.
	
		% *PREPARATORYREDEMPTION
	This also answers a question about Al 13:3, what a preparatory redemption is. It's the new birth. It's a \textit{preparatory} redemption because it isn't the big one that will happen when we are resurrected and glorified; but it's still a redemption, a passing from death unto life through a new birth. It's preparatory to the real one, in the same way that temple ordinances are preparatory to actual becoming a king/queen and priest/priestess. The new birth itself is a preparatory ``ordinance,'' so to speak (and maybe an actual one), that will be fulfilled and actualized when we are resurrected and sanctified.
		See also Al 22:13, where the plan of redemption gets us out of the carnal state.
		See also:
			DC 84:26--27 (preparatory gospel is the gospel of repentance and baptism, and remission of sins); 
		
\item \textbf{Alma's song in Alma 29; compare to the other ``songs'' in the scriptures, like Nephi and the others in the OT}

	Compare also to Nephi's song in Hel 7:7--9.

\item \textbf{Compare what Korihor taught to what the Zoramites believed; why would they have run him through and killed him? Compare what Alma et al. taught in response (Al 30--31; 32--34)}

	For the first part, see 30:13 and 31:22, for example.
	
	For the second part, compare what Alma teaches about the effects of the word being ``real'' because they are ``light'' (32:35), to what Korihor preaches in 30:16.
	
\item \textbf{Discuss how one has to LEARN to live a spiritual life, as told by Alma to Helaman in Al 37:32--37}

	It may not come naturally---someone has to teach us how to do this.
	
\item \textbf{One way the Nephites fall: slowly setting aside what they know until they are left to their own strength (Hel 4)} 

	I read this chapter today, 11/30/2015, and felt a definite rebuke. Lately I haven't been very spiritually diligent, and I think this chapter explains why. I just let things slide and got too casual, relied on myself too much. A good chapter to use to call people to repentance.
	
\item \textbf{The people arrive at the edge of the ocean and the brother of Jared delays four years in calling on the Lord (Eth 2)} 

	Why did he wait four years? I bet he knew what was going to be required of him, and it was too much for him to think about; they were too happy there; his kids were growing up; they had already left behind other places that would have been good to live in; they get to the sea, and he's got to realize what the next step is. But he's scared and comfortable so he doesn't do anything about it. 
	
	What's the next step for us spiritually? What sea are we facing down right now? What are we delaying or avoiding because we already know what's required of us, and how difficult it will be? The Lord didn't let them stop in other good places (Eth 2:7) because he had something better for them; what does the Lord have which is better for us, that we aren't receiving because we aren't calling on him to ask for his help in overcoming the last great hurdle? 
	
\item \textbf{The building of the barges to cross the sea (Eth 2)}

	The brother of Jared asks how they are going to \textit{steer} (19) without light. But they're inside the boat, and it's enclosed, so how could they steer whether they have light inside or not? By ``light'' he must mean, how are we going to know where we're going? This makes the story more interesting and the parallel to us more powerful---sometimes we're going to have to board barges in faith without being able to see where we're going at all. No light illuminates the forward path; all we can do is believe that we are being guided by the Lord in the proper way.
	
	Think of this story in comparison to Nephi's story of crossing the waters---the entire thing, from preparing the boat (Nephi had to learn how to build it/brother of Jared used his old boats as a model), to steering (Nephi had to steer vs. enclosed ship; see also 1N 18:13 vs. Eth 2:19), and so on.
	
	Notice the difference in the questions he asks about light---in 19 it's about steering, while in 22 it's about having light inside the vessel so they don't have to cross in darkness.
	
	Follow up on this with the information that it's 6:1--10.
	
\item \textbf{Use the preparing of the stones as a metaphor to describe how the Lord answers questions, or how he gives a testimony, or how human writings become divine scripture---he takes what is already there and makes it shine with inspiration and divine light and intelligence (Eth 2--3)}

	In our lives we are sometimes swallowed up in the depths of the sea (2:25). But the scriptures are a source of light for us, prepared to illuminate us when nothing else will. How do human writings become inspired? Humans take of the material that they already have, like stones, and think and refine and try to understand things as best they can. Then they approach the Lord with their questions and thoughts, each a precious stone that they have worked on refining. The Lord touches each one, causing each one to bring forth light and have power to guide on its own----he corrects mistakes; he sends a testimony; he causes ideas to shine with power that they didn't have before, by adding intelligence through the spirit; and so on. We take what we have to the Lord and he makes it shine.
	
\item \textbf{The Lord appears to the brother of Jared (Eth 3)}

	A few points to note:
	
	\begin{compactitem}
	\item In 6 it doesn't say that the Lord appeared or arrived; he was already there, talking to the brother of Jared in person. But because of the veil he couldn't see him until a certain moment. Extremely powerful---how close is God to us when we pray? Are there times that he's right next to us, and we just can't discern him? He just stretched forth his hand, already being right there, and the the boj's perception changed. 
	\item The pre-earth Christ, not having received a body, is able to physically interact with the stones. Not immaterial.
	\item The veil is taken off his \textit{eyes}, not his mind or anything else. There's something about the perceptual faculty itself, beginning with the eye, that needs to be changed in order for us to perceive God. The physical matter and processes have to be changed, quickened maybe, in order for us to perceive the Lord. Note verses 19--20 here as well; it may be that the boj had gained so much knowledge, had so much experience with the Lord, that the veil had already been effectively removed from his mind, and what was blocking him know was the physical matter and processes of his perceptual capacities. The mental or intelligent component was already removed from the veil because he had already worked with the Lord so much.
	\item Note also that right before this the boj had been chastised by the Lord for not calling on him for four years (!). Then soon after that chastising experience, the Lord asks him how he wants to cross the waters with light (showing confidence in him), and then revealing himself to him (manifesting his worthiness). Extremely important I think.
	\end{compactitem} 
	
\end{itemize}

\subsection{1 Nephi} *1N

	\subsubsection{1 Nephi 1} % *1N1

		\begin{compactdesc}

		\item[4--7$\vert$]
		\begin{compactdesc}
		\item[God using what is already there to teach us] The first spiritual experience Lehi has after praying ``in behalf of his people'' is fire coming and dwelling upon a rock. Either from out of the fire, or from a different source, he hears and sees many things, and they scare him. This experience is similar to the one Moses has in Exodus 3. Moses is watching the flocks and takes them a ways away when an angel appears in a bush. The bush burns, but isn't consumed. I assume that Lehi was familiar with this story, so he would probably have recognized the fire as a manifestation of God or the Spirit, and he would have known to pay attention. We know that he was also away from home, perhaps working, because he prays ``as he went forth'' and after the experience he must ``return to his own house at Jerusalem.'' The point here is that God spoke to Lehi in circumstances which were similar to a story he would have been familiar with, maybe to make it easier for Lehi to realize he was having a spiritual experience and easier for him to pay attention. The Lord worked with circumstances and symbols that he was already familiar with. He couldn't grow unless he recognized that he was experiencing and opportunity for growth.
		\end{compactdesc}

		\end{compactdesc}

	\subsubsection{1 Nephi 2} % *1N2

		\begin{compactdesc}

		\item[20$\vert$]
		\begin{compactdesc}
		\item[A land of promise...which I have prepared for you] I'm interested in how their understanding and perception of this idea changed over time. Just to note a couple ideas briefly---Jacob later says that there is a happiness \textit{which is prepared} for the saints (2N 9:43). But later the same Jacob talks about how lonely and unhappy everybody was, at least during certain times of their lives; Jac 7:26 is the best example of this, especially at the end, where he says that the people \textit{did mourn out [their] days}. So I'm wondering if, in particular, Jacob's view of the ``land of promise'' changed over time: he thought it was actually the land they started to live on after traveling in the boat, but as the people became less happy, Jacob may have come to think of that ``promised land'' which was prepared as more of a metaphor for the kind of blessed state that the Lord would give to the righteous and the saints. I haven't checked in other places for this yet, but I'm struck by Jac 7:26 in particular. There are more things going on, though, as Jac 7:23 makes things sound a little better.
		\end{compactdesc}

		\end{compactdesc}
		
	\subsubsection{1 Nephi 3} % *1N3
	
		\begin{compactdesc}
		
		\item[7$\vert$]
		\begin{compactdesc}
		\item[save he shall prepare a way] See CDBY 38: ``When God sends a man to do a work, all the devils in hell cannot kill him until he gets through his work; so with Joseph, he prepared all things, gave the keys to men on the earth, and said, ``I may soon be taken from you.' ''
			See also:
				\begin{compactdesc}
				\item[1N 17:2--3] And so great were the blessings of the Lord upon us, that while we did live upon raw meat in the wilderness, our women did give plenty of suck for their children, and were strong, yea, even like unto the men; and they began to bear their journeyings without murmurings. 3 And thus we see that the commandments of God must be fulfilled. And if it so be that the children of men keep the commandments of God he doth nourish them, and strengthen them, and provide means whereby they can accomplish the thing which he has commanded them; wherefore, he did provide means for us while we did sojourn in the wilderness.
				\item[1N 17:13] And I will also be your light in the wilderness; and I will prepare the way before you, if it so be that ye shall keep my commandments; wherefore, inasmuch as ye shall keep my commandments ye shall be led towards the promised land; and ye shall know that it is by me that ye are led.
				\end{compactdesc}
		\end{compactdesc}
		
		\end{compactdesc}

	\subsubsection{1 Nephi 4} % *1N4

		\begin{compactdesc}

		\item[1$\vert$]
		\begin{compactdesc}
		\item[Let...for behold...therefore (verse 2)] Note that the structure of this verse, and the \textit{therefore} that begins verse 2, is similar to 3:16 that comes before it. If I ever look into the particles and structure then keeping consistent with authors will help. Nephi uses this construction twice here fairly close together, and there could be more instances that I've missed already.
		\end{compactdesc}

		\item[3$\vert$]
		\begin{compactdesc}
		\item[the Lord is able...to destroy Laban] In the chapter previous, the angel tells Nephi that the Lord would deliver Laban into his hands. Nephi then assures his brothers, in this verse, that the Lord can ``destroy'' Laban if he wishes. But when it comes time for the Lord to do that---through his servant, Nephi---Nephi doesn't want to do it. In 4:11 when the Spirit is speaking to him, it seems as if the Spirit is saying, ``The Lord has delivered him into your hands. What did you think was going to happen? The Lord is able to destroy him---did you think that was going to be by a bolt of lightning or something? The Lord works through his servants; this is what he is calling upon you to do.''

		There is a marvelous and powerful parallel between Nephi's experience here and the Lord's own experience with the atonement as he describes it in DC 19. I think we should look at Nephi's experience here as the kind of sacrifice described in DC 138:13, as a sacrifice made by the faithful who are waiting for Christ. It's something that the faithful have done---they have overcome themselves in order to follow the Spirit and fulfill the will of God, just as Christ did. The sacrifice is tailored to the individual circumstances and dispositions of the person who will make it. We know that Nephi was very concerned with following the word of the Lord; sometimes it seems like he's obsessed wih that. He's obsessed with following the commandments. And his sacrifice is to go against a commandment as it is written, against a prohibition that isn't just a commandment but that is a principle that I'm sure he knows to be good. And he's asked to violate that. 
		\end{compactdesc}

		\item[6$\vert$]
		\begin{compactdesc}
		\item[Led by the Spirit, not knowing beforehand...] This is like the definition of ``living by every word that proceeds forth from the mouth of God'' (DC 84:44). This is where we need to be at in order to offer the kinds of sacrifices that Nephi is describing here, I think. We don't know beforehand what they are. We're asked to do many basic things, and these are sacrifices that everyone is asked to make. But what is \textit{my} sacrifice? What am \textit{I} supposed to offer up on the altar? I can't know that beforehand. I have to learn it, and to learn it I have to live by the Spirit and by every word that proceeds forth from the mouth of God.
		\end{compactdesc}

		\item[10$\vert$]
		\begin{compactdesc}
		\item[And I shrunk and would that I might not slay him] The parallel to Christ and to great sacrifices made by those who follow him is strengthened in these words, I think---the compare very directly to DC 19:18, where the Lord is describing the atonement and how the thought of going through with it made him feel: ``[I] would that I might not drink the bitter cup, and shrink.''
		\end{compactdesc}

		\item[16--17$\vert$]
		\begin{compactdesc}
		\item[I also knew.../I knew...] Maybe translate this ``know'' as \textit{realize}. Nephi \textit{realized} that these things. This is a moment of great spiritual clarity for him, when everything falls into place and he realizes that this is what he must do. I used to view these verses as Nephi's rationalization of what he was being asked to do; like he's trying to go through the reasons in favor it, trying to convince himself to do it. I don't view them that way anymore. Moments of spiritual clarity like this are where all the pieces fall into place for us, where past thoughts and present questions and future directions are united in a moment of clear understanding. That is what Nephi had here. He suddenly understood all these things at once, and he saw how everything fit together. Once he saw that, there was no hesitation; he obeyed right away. This is such a powerful story in so many ways.
		
		For other moments of undeniable clarity like this one, where events from the past connect with what we are presently experiencing and everything makes sense, see Hel 9:4--5 (plus the verses around for context);
		\end{compactdesc}
		
		\item[18$\vert$]
		\begin{compactdesc}
		\item[own sword/cut off head] See 1Sam 17:51, where David does the same thing to Goliath.
		\end{compactdesc}

		\item[35$\vert$]
		\begin{compactdesc}
		\item[Zoram] The Hebrew root \he{zrm} means \textit{pour forth in floods, flood away}, and there is a noun, \he{zErEm}, which means \textit{flood of rain, rainstorm, downpour}. I can't think of how this would relate to Zoram at all, or how it might relate to his descendants, at least right now.
		\end{compactdesc}

		\end{compactdesc}
		
	\subsubsection{1 Nephi 5} % *1N5
	
		\begin{compactdesc}
		
		\item[5$\vert$]
		\begin{compactdesc}
		\item[I have obtained a land of promise] Weird tense. See 2N 1:3, where maybe Lehi is speaking about L/L and his kids in particular, not including himself. Note the different pronouns in that verse. Then see verse 5 in that chapter as well.
		\end{compactdesc}
		
		\end{compactdesc}
		
	\subsubsection{1 Nephi 7} % *1N7
	
		\begin{compactdesc}
		
		\item[1$\vert$]
		\begin{compactdesc}
		\item[was not meet to be alone] Maybe see the entire Lehi/Nephi story as a type of our lives and our death, and our entering into the celestial kingdom? Sacrifice and covenants required; you can't enter into the highest degree without being married; stuff like that. Just one way to look at everything maybe.
		\end{compactdesc}
		
		\end{compactdesc}

	\subsubsection{1 Nephi 8} % *1N8

		\begin{compactdesc}

		\item[1$\vert$]
		\begin{compactdesc}
		\item[Fit in context?] This verse just seems out of place. It doesn't seem related at all to what comes before---a brief report of how they reacted to Nephi and the others returning with Ishmael's family. It doesn't seem related to what follows, either, which is the beginning of Lehi's vision of the tree of life. I thought that maybe it just popped into Nephi's head as he is about to report the vision, because this verse involves fruit and so does the vision, but that also seems a bit odd. I'm not sure what to make of it here.
		\end{compactdesc}
		
		\item[30$\vert$]
		\begin{compactdesc}
		\item[they did press their way forward...] In our own life, we will repeat this rod/temptation/fruit cycle many times; it's like a videogame sometimes, pressing the reset button, and we get put back at the start of the rod to try again with different enemies on a different level, maybe a different level of difficulty in some ways. Also, note that we are traversing multiple rods at the same time in different aspects of our life. In some aspects we are eating the fruit of the love of God right now, in others we are struggling through the mists, and maybe in some others we're headed to the great and spacious building or are already there. All we can do is try to move forward on all the different paths and do our best to stay at the tree once we arrive---and when we inevitably leave the tree, we just try to get back their again.
		\end{compactdesc}

		\item[30--34$\vert$]
		\begin{compactdesc}
		\item[he...he...his/me...we] Notice the shift in pronouns from verses 30--32 to 33. It may be that the beginning of verse 34, \textit{these are the words of my father}, is supposed to apply to 33 and perhaps the remainder of 34 as well. TET for 34 is different: ``Thus is the words of my father, \textbar for as many as heeded them had fallen away.'' That makes it seem even more likely that the beginning of 34 applies to what comes before it, not what comes after it; what comes after it is kind of an explanation or a summary of what's in 33. Editing not very good for the 1981 edition here. 

		There are other instances of things like this in the Book of Mormon, I think, with shifting pronoun reference. I think I've seen it before with some of the stuff that Mormon says when he reports other stories, but I can't remember.
		\end{compactdesc}

		\end{compactdesc}

	\subsubsection{1 Nephi 10} % *1N10

		I'm struck by what Nephi says here in verse 1: \textit{And now I, Nephi, proceed to give an account upon these plates of my proceedings, and my reign and ministry.} So what was he giving before? An account of his father's ``reign and ministry''? I don't think so. I think we can view the first nine chapters as introductory material which sets the stage for the story that Nephi intends to tell about himself. This includes the account of the family leaving Jerusalem, the two return trips, and Lehi's vision. Nephi then says he is going to write about his own reign and ministry, and finishes the thought with, \textit{wherefore, to proceed with mine account, I must speak somewhat of the things of my father, and also of my brethren.} In other words, to make sure the stage is fully set, he has to say a few more things about what his father has been preaching---some prophecies about the Messiah, John the Baptist, and the olive tree.

		Once he has reported these things, everything is set. He finishes all of this introductory material with some brief comments about the Holy Ghost and how knowledge is available to everyone by its power. One thing I think he is saying here is that the kind of revelatory experience that he is going to report next is open to everyone.

		So what is the turning point between the introductory, stage-setting material, and the point where Nephi takes over the record with his own reign and ministry? I think it's the point where he has his own revelatory, mind-opening experience in the next few chapters. In 10:17 he says that he is ``desirous also that I might see, and hear, and know of these things,'' which echoes what he says in 2:16. But in that earlier occasion, his heart is softened to believe the words of his father--it doesn't look like he receives any additional light or knowledge. He receives a testimony of the things which are being told to him, but that is all. However, on this later occasion, he receives new things and a greater testimony of his own. This is the important turning point for him, I think. He came into his own spiritually. He had a great trial of faith, a test, in going back to get the plates. Once he passed the test, he was ready to receive more. With all of those stories written down, Nephi is ready to tell about himself and his spiritual experiences. This also gives me a new perspective on Eth 12:6. Maybe we can read it with a small addition and some emphasis: \textit{wherefore, dispute not because ye [do not] see [yet], for ye receive no [personal] witness until after the trial of \textbf{your} faith.} Until that time of a personal witness, we may have to be content with a testimony of what someone else has said, or just trust in their words.

		\begin{compactdesc}

		\item[17--22$\vert$]
		\begin{compactdesc}
		\item[Retrospective?] It also strikes me that this set of verses definitely look like something that Nephi is saying as he looks back on his life. Obviously this entire story comes from Nephi's retrospective, but I think these verses in particular are not something that Nephi would have written in the early period of his life, as they were still in the desert. I think he's saying this with the experience of years, and trying to encourage us to look for the Lord and greater knowledge through the Holy Ghost in our own lives, wherever or whenever we live.

		Also see the note on \textit{deny} in verse 22.
		\end{compactdesc}

		\item[23$\vert$]
		\begin{compactdesc}
		\item[deny] This is the first use of \textit{deny} in the BOM, and we don't get another until 2N 25. There are 10 uses in 2N.
		\end{compactdesc}

		\end{compactdesc}
		
	\subsubsection{1 Nephi 11} % *1N11
	
		\begin{compactdesc}
		
		\item[18$\vert$]
		\begin{compactdesc}
		\item[after the manner of the flesh] Extremely important qualifier on the mother of Christ. Not after the manner of the spirit, or whatever else you want to say.
		\end{compactdesc}
		
		\end{compactdesc}

	\subsubsection{1 Nephi 16} % *1N16

		\begin{compactdesc}

		\item[23$\vert$]
		\begin{compactdesc}
		\item[Asking Lehi where to go] I think I've heard in the past (from President Openshaw, maybe) that Nephi asks Lehi where to go here because Nephi recognizes his father's role as the leader of the expedition and the one who could properly receive the revelation about it. But it may also be that Nephi takes the chance to ask him in order to prod his father to repent and get back on the wagon. In verse 22 Nephi speaks unto his \textit{brethren}, which could include his father, but I would have to look into how that word is used. What I'm thinking is that he spoke in 22 to his brothers and their families, and then asked his father about where to go in order to remind his father that he needed to quit complaining and solve the problem the right way.

		What's interesting is that the directions Nephi then follows, in verse 30, are given from the ball, not from Lehi.
		\end{compactdesc}

		\item[23$\vert$]
		\begin{compactdesc}
		\item[beheld] \textit{Behold} here (in the past tense in this verse) has to mean something like \textit{understand} or \textit{come to know}, \textit{realize}. He understands for the first time how the ball works.
		\end{compactdesc}

		\end{compactdesc}

	\subsubsection{1 Nephi 17} % *1N17

		\begin{compactdesc}
		
		\item[14$\vert$]
		\begin{compactdesc}
		\item[ye shall know that I, the Lord, am God] What kind of revelation is this? Something to do with the temple in the new world? I'm wondering how the Lord would show them this, or what it would be like to learn.
		\end{compactdesc}

		\item[21$\vert$]
		\begin{compactdesc}
		\item[and we might have been happy] In my opinion, one of the saddest things uttered in the Book of Mormon, both because it reveals how miserable Laman and Lemuel must have been during most of the trip from Jerusalem to Bountiful, and also for another reason. It shows how little Laman and Lemuel changed as a result of all their experiences; there was almost no spiritual growth for them, almost nothing to show for all those years. Very sad.

		So what was the difference? I remember on my very first real day in the mission, September 28, President Openshaw or the assistants taught a very similar lesson--they pointed out that Nephi and Laman/Lemuel had done many of the same things, but the results were different. Part of this is recorded in my mission journal. Here's the little table they drew up:

		\begin{tabular}{|c|c|c|}
		\hline Nefi \& Mandamientos \& Laman y Lemeul \\ 
		\hline \ding{51} \& salir al desierto \& \ding{51} \\ 
		\hline \ding{51} \& buscar las planchas \& \ding{51} \\ 
		\hline \ding{51} \& buscar esposas \& \ding{51} \\ 
		\hline \ding{51} \& construir barco \& \ding{51} \\ 
		\hline deseo \& \& deb\'{\i}an hacerlo \\ 
		\hline 
		\end{tabular}

		Laman and Lemuel ended up participating in all the big things that Nephi did, but they were never changed. They weren't the fourth missionary. They did not spiritually benefit from any of the things they did. 

		So what made the difference for Nephi? I can think of two things:

		\begin{enumerate}
		\item He sought testimony, knowledge, and inspiration directly from the Lord.
		\item He strove to obey the commandments cheerfully and quickly.
		\end{enumerate}

		Was there really any other difference? I'm not sure. I don't think it needs to be more subtle than that. Nephi did those two things, and Laman and Lemuel didn't. Nephi's heart changed, and Laman's and Lemuel's didn't. 
		\end{compactdesc}

		\item[25--29$\vert$]
		\begin{compactdesc}
		\item[Parallels to Lehi's family] So in these verses Nephi gives a list of a number of things that Laman and Lemuel \textit{know}, or that at least Nephi claims they know. I think that there are some possible parallels between the events that Nephi mentions, and the things that have happened to Lehi's family until this point. Here they are:

		\begin{enumerate}
		\item \textbf{ye know that the children of Israel were in bondage}
		\item \textbf{ye know that they were laden with tasks, which were grievous to be borne}
		\item \textbf{ye know that it must needs be a good thing for them, that they should be brought out of bondage}
		\item \textbf{ye know that Moses was commanded of the Lord to do that great work}
		\item \textbf{ye know that by his word the waters of the Red Sea were divided hither and thither, and they passed through on dry ground}
		\item \textbf{ye know that the Egyptians were drowned in the Red Sea, who were the armies of Pharaoh}
		\item \textbf{ye also know that they were fed with manna in the wilderness}
		\item \textbf{ye also know that Moses, by his word according to the power of God which was in him, smote the rock, and there came forth water, that the children of Israel might quench their thirst}
		\end{enumerate}

		How legitimate are these parallels?
		\end{compactdesc}

		\item[25$\vert$]
		\begin{compactdesc}
		\item[ye know...] There is an interesting connection here between these verses and an earlier comment that Nephi makes. Here Nephi is trying to reason with his brothers, and he brings up the story of Moses and the children of Israel. He uses the words ``you know'' no less than 8 times in this verse and the next four verses. At the beginning of 1 N 4, Nephi says,

		\begin{quote}
		let us be strong like unto Moses; for he truly spake unto the waters of the Red Sea and they divided hither and thither, and our fathers came through, out of captivity, on dry ground, and the armies of Pharaoh did follow and were drowned in the waters of the Red Sea. Now behold ye know that this is true; and ye also know that an angel hath spoken unto you (2-3).
		\end{quote}

		In both places, when trying to reason with his brothers, Nephi brings up the story of Moses and the relationship that his brothers have to the story: they \textit{know} it. The story seems to have a special place for them. Perhaps it is a story that they enjoyed when they were younger, or that they pretended to act out when they were playing. Or maybe they had some other connection to it that made it special for them. I'm not sure. But Nephi brings it up for a reason, and Laman and Lemuel don't ever deny that they know it, at least as far as I can tell. 
		\end{compactdesc}

		\item[30$\vert$]
		\begin{compactdesc}
		\item[you know] The \textit{you know} language from the past five verses stops here, and what we get in its place is a \textit{notwithstanding}. The parallel he is trying to draw is obvious--they have already experienced many miracles themselves, just like the children of Israel on their way out of Egypt, and despite that they are hardening their hearts and not listening to the Lord's commandments, just like the children of Israel. So he talks a great deal about what they \textit{know}, and then he goes on to reason from that basis toward why they need to change their attitudes.
		\end{compactdesc}

		\item[31$\vert$]
		\begin{compactdesc}
		\item[there was not any thing done save it were by his word] The way that this sentence is set up reminds me of 2 N 9:20: ``there is not anything save he [God] knows it.'' See also 3 N 8:1 for a similar construction, as well as Morm 8:10, 36.

		Beyond that, though, I'm not really sure what to make of this part of the verse. 
		\end{compactdesc}

		\end{compactdesc}
		
	\subsubsection{1 Nephi 18} % *1N18
	
		\begin{compactdesc}
		
		\item[23$\vert$]
		\begin{compactdesc}
		\item[and we did call it the promised land] Kind of an interesting comment. The promised land is where we make it; it's where we end up after we've been working with the Lord. Reminds me of 1N 5:5, ``I have obtained a land of promise.'' That was long before they got into the ship.
		\end{compactdesc}
		
		\end{compactdesc}
		
	\subsubsection{1 Nephi 19} % *1N19
	
		\begin{compactdesc}
		
		\item[2--6$\vert$]
		\begin{compactdesc}
		\item[the plates] Note that the earliest plates Nephi makes come long after they've left Jerusalem, so the record from those days is faded and just part of Nephi's memory.
		\end{compactdesc}
		
		\end{compactdesc}

\subsection{2 Nephi} % *2N

	\subsubsection{2 Nephi 2} % *2N2

		\begin{compactdesc}

		\item[5$\vert$]
		\begin{compactdesc}
		\item[men are instructed sufficiently that they know good from evil] 2N 25:28.

		Here are the requirements to be a moral agent that Dan Dennett give in his lecture ``The Moral Agents Club'' on September 21, 2012. He called them ``specs for a morally competent agent'': is well informed, has roughly well-ordered desires, is moved by reasons (can see reasons in the world and can be moved by verbally presented reasons --that we can reason with it), it ``could have done otherwise,'' is not being controlled by another agent, is punishable (has to have skin in the game; we must be able to punish it, it needs to be vulnerable. Mutual vulnerability is a key ingredient).
		\end{compactdesc}
		
		\item[7$\vert$]
		\begin{compactdesc}
		\item[answer] Definition \#3 from 1828: ``To comply with, fulfill, pay or satisfy''; definition \#6, ``To perform what was intended; to accomplish.''
		\end{compactdesc}
		
		\item[13$\vert$]
		\begin{compactdesc}
		\item[And if there is no God we are not...all things must have vanished away] This is a strange thing to say because we know it has to be false, based on what JS later taught about intelligences and matter being self-existent (they exist independently of God). I guess you could look at it as Lehi getting it wrong, but I like to see it as an indication of JS's thought developing here. This is very early, mid-1829 when we get this portion of the BOM. His thinking progressed on this issue. So I think we have to reject the most natural interpretation and interpret it a different way.
		\end{compactdesc}
		
		\item[14$\vert$]
		\begin{compactdesc}
		\item[he hath created all things...things to act and things to be acted upon] Not sure how to take this either. We know it's false when you take it on the most natural interpretation. We have to understand it some other way.
		\end{compactdesc}

		\end{compactdesc}

	\subsubsection{2 Nephi 9} % *2N9

		\begin{compactdesc}

		\item[7$\vert$]
		\begin{compactdesc}
		\item[atonement] The atonement needs to be infinite in order for corruption to put on incorruption. I believe that this means the atonement must be infinite for me--that if Christ had only performed the atonement for me, it would still have had to be infinite, because this is what is required for any individual to go from corruption to incorruption. It is a marvelous thought, that the atonement is infinite for me. See also 2 Nephi 25:16, ``the atonement, which is infinite for all mankind.'' Note that 2 Nephi 9:7 contains the first and two total of the four occurrences of the idea of an ``infinite'' atonement.
		\end{compactdesc}
		
		\item[41$\vert$]
		\begin{compactdesc}
		\item[the keeper of the gate...none other way save it be by the gate] See 2N 31:17, where the ``gate'' is repentance and baptism by water.
		\end{compactdesc}

		\end{compactdesc}

	\subsubsection{2 Nephi 10} % *2N10

		\begin{compactdesc}

		\item[24$\vert$]
		\begin{compactdesc}
		\item[reconcile] It is also only through the grace of God that we are reconciled to the will of God; see Jacob 4:11.
		\end{compactdesc}

		\end{compactdesc}

	\subsubsection{2 Nephi 12} % *2N12

		\begin{compactdesc}

		\item[3$\vert$]
		\begin{compactdesc}
		\item[he will teach us of his ways] See DC 84:48.
		\end{compactdesc}
		
		\item[16$\vert$]
		\begin{compactdesc}
		\item[ships of the sea/Tarshish/pleasant pictures] See note at Isa 2:16.
		\end{compactdesc}

		\end{compactdesc}

	\subsubsection{2 Nephi 13} % *2N13

		\begin{compactdesc}

		\item[9$\vert$]
		\begin{compactdesc}
		\item[show of their countenance doth witness against them] See Hel 9:33--34.
		\end{compactdesc}

		\item[10$\vert$]
		\begin{compactdesc}
		\item[say unto the righteous that it is well with them] See 2N 28:7--8.
		\end{compactdesc}

		\item[12$\vert$]
		\begin{compactdesc}
		\item[they who lead thee cause thee to err and destroy the way of thy paths] See Alma 46:9, Mosiah 11:4-7.
		\end{compactdesc}

		\end{compactdesc}

	\subsubsection{2 Nephi 15} % *2N15

		\begin{compactdesc}

		\item[4$\vert$]
		\begin{compactdesc}
		\item[What could have been done more to my vineyard that I have not done in it?] See Jacob 5:47.
		\end{compactdesc}

		\item[13$\vert$]
		\begin{compactdesc}
		\item[my people are gone into captivity, because they have no knowledge] See Al 12:11.

		For the Hebrew word ``da'at'' which is translated as ``knowledge'' here, see BDB 395b. See also the LXX here: ``Indeed my people have become captive because they do not know (about [\gr{eidenai}]) the Lord.''
		\end{compactdesc}

		\end{compactdesc}

	\subsubsection{2 Nephi 16} % *2N16

		\begin{compactdesc}

		\item[7$\vert$]
		\begin{compactdesc}
		\item[thine iniquity is taken away] Note that ``iniquity'' is singular; cf. 3 Nephi 8:1, ``there was not any man who could do a miracle in the name of Jesus save he were cleansed every whit from his iniquity,'' also singular.

		\item[...and thy sin purged] See Hebrews 1:3, which notes that Christ is the one who has done this.
		\end{compactdesc}

		\end{compactdesc}

	\subsubsection{2 Nephi 21} % *2N21

		\begin{compactdesc}

		\item[9$\vert$]
		\begin{compactdesc}
		\item[knowledge] ``Knowledge'' in this verse is not ``da'at'' but ``da'eh,'' a form related to 2 Nephi 15:13 but distinct from it. Note that here, as in the LXX in 2 Nephi 15:13, it is explicit that the knowledge is of ``the Lord.'' In this verse, the LXX has \textit{tou gnonai} (articular infinitive, from \textit{ginosko}) instead of \textit{to eidenai}. For the sake of comparison, \textit{ginosko} appears 746 times in the LXX, while \textit{oida} (whence \textit{eidenai}) appears 283 times.
		\end{compactdesc}

		\end{compactdesc}

	\subsubsection{2 Nephi 22} % *2N22

		\begin{compactdesc}

		\item[3$\vert$]
		\begin{compactdesc}
		\item[Therefore, with joy shall ye draw water out of the wells of salvation] See John 4:13-14. Note that in John 4:14 the water shall ``be'' in us a well, but in Greek it's actually ``become.'' The Greek word ``springing up'' is \textit{hallomai}, which means to ``well'' or ``well up,'' which in turn means ``to rise up to the surface and spill or be about to spill.'' I think we often miss that the word denotes overflowing or abundance.

		See also 2 Nephi 21:9, where the knowledge of the Lord fills the earth in a manner ``as the waters cover the sea,'' again in great abundance. It is literally overflowing.
		\end{compactdesc}

		\end{compactdesc}

	\subsubsection{2 Nephi 24} % *2N24

		\begin{compactdesc}

		\item[1$\vert$]
		\begin{compactdesc}
		\item[and will yet choose Israel] ``Choose'' here is \textit{bahar} (\textit{het} as the second radical; the object is marked with \textit{b-}). It is pretty common and means just ``to choose.'' The overwhelming majority of instances of the verb ``choose'' in the OT are this word.
		\end{compactdesc}

		\item[16$\vert$]
		\begin{compactdesc}
		\item[Lucifer] Lucifer is called an \textit{ish} in this verse, a ``man.''
		\end{compactdesc}

		\item[24$\vert$]
		\begin{compactdesc}
		\item[coming to pass] Not sure how I ought to understand the part about the Lord thinking and something coming to pass.
		\end{compactdesc}

		\end{compactdesc}

	\subsubsection{2 Nephi 25} % *2N25

		\begin{compactdesc}

		\item[6$\vert$]
		\begin{compactdesc}
		\item[unto my children] The first ``unto my children'' in this verse is not in Skousen's earliest text of the BOM.
		\end{compactdesc}

		\item[12$\vert$]
		\begin{compactdesc}
		\item[When the day cometh that the Only Begotten of the Father...shall manifest himself unto them in the flesh] Nephi here is talking about the Jews, so he is referring to the time when Christ would come into the world and be among those people. Compare this with 2 Nephi 32:6: ``And when he shall manifest himself unto you in the flesh, the things which he shall say unto you shall ye observe to do.'' The question is, who is the ``you'' to whom Nephi addresses himself in 2 Nephi 32:6? In 2 Nephi 25:3 he says that he writes ``unto my people, unto all those that shall receive hereafter these things which I write.'' I believe he is directing himself to the same group of people (us, we who receive the things which he has written) in 2 Nephi 32:6. This matters because one interpretation of 32:6 is that Christ will manifest himself in the flesh to us as part of our spiritual progression (DC 88:68, John 14:23/DC 130:3), rather than that this verse refers to his apperance among the Nephites beginning in 3 Nephi 11. See also 2 Nephi 26:13, where ``in the flesh'' is conspicuously absent. See also 2 Nephi 26:1.
		\end{compactdesc}

		\item[13$\vert$]
		\begin{compactdesc}
		\item[with healing in his wings] See Malachi 4:2 (3:20 in BHS/LXX), ``the Sun of righteousness arise with healing in his wings.'' The pun only works in English, as the BHS has \textit{shemesh} and the LXX has \textit{helios}.
		\end{compactdesc}

		\item[14$\vert$]
		\begin{compactdesc}
		\item[After the Messiah hath risen from the dead...] See 2 Nephi 26:13, ``he manifestesth himself unto all those who believe in him, by the power of the Holy Ghost,'' and 3 Nephi 16:4, ``But now I go unto the Father, and also to show myself unto the lost tribes of Israel, for they are not lost unto the Father, for he knoweth whither he hath taken them.''
		\end{compactdesc}

		\item[16$\vert$]
		\begin{compactdesc}
		\item[infinite atonement] This is the third of four occurrences of the idea of the ``infinite atonement.'' All four occurrences are in the BOM. The next and last is probably the most well-known, Alma 34:12.

		\item[pure hearts and clean hands] The only other two instances are Psalms 24:4 and Alma 5:19. Psalms 24:3 looks to contain a temple reference.
		\end{compactdesc}

		\item[19$\vert$]
		\begin{compactdesc}
		\item[Jesus] This verse is the first occurrence of ``Jesus'' in the BOM and so the first occurrence of the full name ``Jesus Christ.''

		\item[And according to the words of the prophets...] See 2N 10:3.
		\end{compactdesc}

		\item[20$\vert$]
		\begin{compactdesc}
		\item[poisonous serpents] The story of Moses poisonous serpents and healing by the raised serpent appears in at least five distinct places in the BOM: 1 Nephi 17:41, 2 Nephi 25:20, Alma 33:19, Alma 37:46, Helaman 8:14. The only other OT stories that I can think may be more commonly referenced than this one are the Adam and Eve story and the story of the Israelites escaping Egypt and crossing the Red Sea.
		\end{compactdesc}

		\item[22$\vert$]
		\begin{compactdesc}
		\item[shall be judged of them] ``And the nations who shall possess them [these things; the things that Nephi writes] shall be judged of them according to the words which are written'' --I think that we can generalize this statement to all people and nations, that they shall be judged according to the word of God which is written among them and which they actually have reasonable access to. It is a very important principle, especially when one considers that the ``word of God written among them'' may not have to be canonical LDS scripture but may instead be influential philosophical or wisdom literature which forms a part of the general morals and culture of a society. In many cases those ideas harmonize with principles of the restoration gospel, and I believe that people will be held accountable for their adherence to them. See also Words of Mormon 1:11.
		\end{compactdesc}

		\item[23$\vert$]
		\begin{compactdesc}
		
		\item[we labor diligently to write...] The exact phrase ``reconciled to God'' appears two other times in the scriptures, in Romans 5:10 and 2 Cor 5:20. In both cases the verb is \textit{katalasso} (see GELNT 521a, b). The importance of this word and idea in the NT cannot be overstated, especially since the noun form is the word used in the only instance of the English word ``atonement'' in the NT (Romans 5:11). The definition for the verb is powerfully descriptive: ``the exchange of hostility for a friendly relationship, reconcile.''

		Note that 2 Nephi 10:24 and 2 Nephi 25:23 both contain the idea of reconciling ourselves to God while simultaneously affirming the utter necessity of God's grace in our salvation.

		Apparently (as I infer from 2 Nephi 10:24), the reconciliation is something we can do ourselves (or perhaps choose to begin and continue doing; certainly it seems that God helps us along the way). I should look to the NT for more about the process of reconciliation.
		
		\item[after all we can do] There are a few ways of reading this phrase. One is the way we normally read it I think, which is, we have to do everything we can do, and \textit{then} we're saved because of grace. Or you can read it a different way: Nephi is saying, look around at everything we try to do for ourselves, all the ways we try to improve ourselves or make ourselves clean or take control of our own lives. After you see all those things that we can try to do for ourselves, it's clear that it's by \textit{grace} we're saved. And there's another way too: ``it's by grace that we are saved, [even] after all we can do [to mess it up].''
		
		\end{compactdesc}

		\item[24$\vert$]
		\begin{compactdesc}
		\item[with steadfastness] ``Steadfastness'' means ``resolutely or dutifully firm and unwavering.'' See Jacob 3:2 and Moroni 7:30.
		\end{compactdesc}

		\item[28$\vert$]
		\begin{compactdesc}
		\item[the words...are sufficient to teach...] See 2 Nephi 2:5, ``and men are instructed sufficiently that they know good from evil,'' and Moroni 7:16, ``the Spirit of Christ is given to every man, that he may know good from evil.'' There are at least two essential components to knowing good from evil (whatever that amounts to): one is a natural endowment (Spirit of Christ), and the other ought to be imparted by others (see also DC 93:40 and 68:25).
		\end{compactdesc}

		\item[29$\vert$]
		\begin{compactdesc}
		\item[the right way is to believe in Christ...] See verse 28; also Mor 7:17.
		\end{compactdesc}

		\end{compactdesc}
		
	\subsubsection{2 Nephi 26} % *2N26
	
		\begin{compactdesc}
		
		\item[13$\vert$]
		\begin{compactdesc}
		\item[manifesteth himself...by the power of the Holy Ghost] See 3N 15:23, ``And they understood me not that I said they shall hear my voice; and they understood me not that the Gentiles should not at any time hear my voice---that I should not manifest myself unto them save it were by the Holy Ghost.'' See also 3N 16:6;
		\end{compactdesc}
		
		\item[24$\vert$]
		\begin{compactdesc}
		\item[he doeth not anything save it be for the benefit of the world] See Eth 3:4, ``And I know, O Lord, that thou hast all power, and can do whatsoever thou wilt for the benefit of man...''
		\end{compactdesc}
		
		\end{compactdesc}
		
	\subsubsection{2 Nephi 31} % *2N31
	
		\begin{compactdesc}
		
		\item[7$\vert$]
		\begin{compactdesc}
		\item[and witnesseth unto the Father...in keeping his commandments] Notice the similarity between this verse about Christ---``witnesseth unto the Father that he would be obedient unto him in keeping his commandments''---and the sacrament prayer, in which we say that we ``witness unto thee, O God, the Eternal Father, that they are willing to...keep his commandments'' (DC 20:77). The similarity in language goes along with the timeframe---2N 31 is translated very near the end of the BOM translation process, around June 21 1829 according to Watson's timeline. DC 20 is listed as a revelation given on 1830, but according to JSP-D1:117--118, there is a lot of evidence to think that portions of the revelation had been received as early as summer 1829. We also know that Oliver Cowdery's version of DC 20, the ``Articles of the Church of Christ,'' was composed in June 1829, and his document contains the same sacrament prayer with the same passage as 2N 31:7 (JSP-D1:372). This verse in 2N 31 also mentions that Christ is humbling himself ``according to the flesh,'' while in the sacrament prayer we're supposed to eat ``in remembrance of the body of [the] Son.''
			I just love finding these connections in JS's thought. Which has priority here, the BOM translation or the revelation on the sacrament? Not sure, but it's all spun together in one.
		\end{compactdesc}
		
		\item[12$\vert$]
		\begin{compactdesc}
		\item[do the things which ye have seen me do] See 3N 27:21, for example; how are we to accomplish this? See the power to know Christ's ministry under ``powers.''
		\end{compactdesc}
		
		\end{compactdesc}
		
	\subsubsection{2 Nephi 33} % *2N33
	
		\begin{compactdesc}
		
		\item[1$\vert$]
		\begin{compactdesc}
		\item[unto the hearts] ``Unto,'' not ``into.'' Yale has ``unto'' also.
			But see 3N 11:3, where the word pierces people to the very center. See other instances of ``pierce'' as well; there are some things that will get into our hearts whether we want them to or not.
				See also:
					DC 29:7; 
		\end{compactdesc}
		
		\end{compactdesc}

\subsection{Jacob} % *JAC

	\subsubsection{Jacob 4} % *JAC4

		\begin{compactdesc}

		\item[13$\vert$]
		\begin{compactdesc}
		\item[validity of concepts in the scriptures] Validity of concepts in the scriptures---President Openshaw---doesn't matter if they (the events in the scriptures) actually happened, they are what we have and I think we have to assume that their preservation was divinely guided. So assume they are true, then learn principles and truths about God based on that assumption. Remember the conference talk--they were given by the Holy Ghost.
		\end{compactdesc}

		\end{compactdesc}

\subsection{Mosiah} % *MOSIAH

	\subsubsection{Mosiah 2} % *MO2

		Compare King Benjamin's address, and the circumstances in which he gave it, to Samuel's address in 1Sam 12.
		
		\begin{compactdesc}
		
		\item[9$\vert$]
		\begin{compactdesc}
		\item[ears/hear, hearts/understand, minds/mysteries] See also 3N 11:5.
		\end{compactdesc}
		
		\end{compactdesc}

	\subsubsection{Mosiah 3} % *MO3

		\begin{compactdesc}

		\item[16$\vert$]
		\begin{compactdesc}
		\item[by nature, they fall] See Eph 2:3. 
		\end{compactdesc}

		\end{compactdesc}

	\subsubsection{Mosiah 4} % *MO4

		\begin{compactdesc}

		\item[12$\vert$]
		\begin{compactdesc}
		\item[or] I'm wondering if the \textit{or} in this verse is an ``in other words'' \textit{or}, or if it's a true disjunctive \textit{or}. The first is explanatory; the second is an actual option between two things. We're told that we'll grow ``in the knowledge of the glory of him that created [us]'' OR ``in the knowledge of that which is just and true.''

		I am inclined to think that this is a disjunctive \textit{or}. I've been reading in 2 Nephi 2, where Lehi remarks that his son Jacob ``hast beheld in [his] youth [God's] glory.'' From DC 93 we know that the glory of God is intelligence. So back in Mosiah 4:12, if it's a true disjunctive \textit{or}, there are two ways we can grow in knowledge as we do the things that Benjamin instructs: we can either grow in the knowledge of the glory of God--or grow in our knowledge of intelligence, light and truth; or we can grow in the knowledge of other things that are just and true. The reason I like to see it as disjunctive is that it broadens the set of things that we can learn when we are righteous. For the first disjunct, we can gain in our spiritual intelligence and our spiritual capacities; for the second, we can gain knowledge about the world that, while it is spiritual in a certain sense, might best be viewed as separate. 

		If it isn't a disjunctive \textit{or}, then we still might say that the second type of truth we can learn from being righteous still technically falls under the knowledge of intelligence, but seeing it as disjunctive helps me to understand that many different kinds of knowledge and truth are going to flow into my life when I am righteous and understand and respect the relationship that I have with God.

		I'm worried that this analysis is just conflating two meanings of \textit{glory}, though. We can just say that one meaning refers to the external splendor which surrounds God, while another meaning relates to whatever DC 93:36 is about. This may be a problem, but I don't know that it threatens my analysis of these verses here; it would be more relevant to looking back at 2 N 2, for example. 
		\end{compactdesc}

		\end{compactdesc}

	\subsubsection{Mosiah 19} % *MO19

		\begin{compactdesc}

		\item[24$\vert$]
		\begin{compactdesc}
		\item[ended the ceremony] Not sure what this means at all. It could refer to many different things that happened while the men were in the wilderness or after they met the other men whom Gideon sent. Here are some pieces of commentary from the Book of Mormon online site:
		\begin{compactitem}
		\item Verse 24 contains the unique phrase ``after they had ended the ceremony...'' It is unclear what this means. Perhaps there was some formal greeting. Perhaps there was a formal ceremony of surrender by the former Noah-men to the Gideon-men. Perhaps it was simply a word that described the exchange of information between the two groups. Another possibility is that Mormon is again conflating events. Perhaps the meeting of the two groups is occurring at approximately the same time as the exchange of oaths noted below, and it is that exchange of oaths that is the ``ceremony.'' Gideon may have sent his men ``secretly'' while the rest of the city was occupied in the ``ceremony'' of formal submission to the Lamanites.

		John Tvedtnes has suggested that the ceremony would have been a purification ceremony whereby the men were returned to ritual cleanliness after their murder of the king (Tvedtnes, John A. ``The Nephite Purification Ceremony.'' In: \textit{The Most Correct Book}. Cornerstone, 1999, pp. 176--7). This suggestion makes sense in terms of Jewish halakic (purity) laws, but does not necessarily fit the Biblically proscribed actions for a murderer. As Tvedtnes points out, the murderer would not be allowed to live under the law of Moses. While an accidental killing allowed for refuge, it would be hard to consider the death of Noah accidental. Nevertheless, there may have been some aspect of the Nephite culture (or at least this branch of Nephite culture) that provided for a ritual cleansing of those who had committed a similar act.

		\item In Mosiah 19:24 we have a strange statement that isn't explained in the text. It's included as if the reader should be fully aware of what was going on. In review, Gideon had sent men into the wilderness to search for king Noah and those that went with him, and they had encountered part of the group as they were about to return to the land of Nephi. The returning group explained to the men of Gideon that they had put king Noah to death and that the wicked priests had escaped farther into the wilderness. Then comes the strange statement: ``And it came to pass that after they had ended the ceremony, that they returned to the land of Nephi, rejoicing, because their wives and their children were not slain...'' (Mosiah 19:24).

		Hugh Nibley asks, What ceremony? From the commentary on ``the Men of Gideon'' (see Mosiah 19:22) we learned that they were a special party or alliance to bear the name even in his absence---the party of Gideon. Remember the party of Noah had left the town, and Gideon had come to catch up with them and punish their leader. Gideon couldn't leave Noah alone---he was on his heels all the time. So they were hostile parties who were opposed to each other. One was the refugees, and the other was the avenging party following them, so they were hostile. They couldn't go back home together until they had settled, smoked the peace pipe, and had the ceremony. You have to have a ceremony before you can reach peace with a hostile group. You either fight them or have the ceremony, so that's what they did. They had a peace ceremony. They always have that, but this is putting it so casually, as if Joseph Smith knew exactly what he was saying. They carried out certain rites of reconciliation here, which is very common and has to be done, as far as that goes. It's unthinkable to omit it. Then they went back and told Gideon himself all that they told the men of Gideon about the king, his old rival. [Hugh W. Nibley, Teachings of the Book of Mormon, Semester 2, p. 125]
		\end{compactitem}

		\end{compactdesc}

		\end{compactdesc}

	\subsubsection{Mosiah 20} % *MO20

		The story in 1--5 about the priests carrying off the dancing girls is similar a story at the end of Judges, 21:16--23.

	\subsubsection{Mosiah 21} % *MO21

		\begin{compactdesc}

		\item[33$\vert$]
		\begin{compactdesc}
		\item[considering himself an unworthy servant] Maybe Ammon realized that something bigger was going on here than he was prepared to handle; Alma, a former priest, had already formed a church, and maybe Ammon thought that something like that needed to happen with Limhi's people as well. So he figured that he wasn't prepared to take on something that big, or wasn't the right person to be in charge of organizing them into a church. The next verse, 34, seems to say that if the people had been baptized at this time, they would have formed a church. Ammon, knowing that Alma was somewhere leading a church of 450 (Mo 18:35), and knowing that Mosiah was back in Zarahemla as the proper king, thought that it wasn't his place to baptize them at this time. See also Mo 23:16--17;
		\end{compactdesc}

		\end{compactdesc}
		
	\subsubsection{Mosiah 29} % *MO29
	
		Compare the issue of a king vs. some other government here to how the people want a king (and the problems with that) in 1Sam 8--12.

\subsection{Alma} % *ALMA

	\subsubsection{Alma 12} % *AL12
	
		\begin{compactdesc}
		
		\item[9--11$\vert$]
		\begin{compactdesc}
		\item[THE PRINCIPLE] \textbf{}
			\begin{compactdesc}
			\item[Related verses] Hel 13:8 (uses ``word'');
			\end{compactdesc}
		\end{compactdesc}
		
		\item[10$\vert$]
		\begin{compactdesc}
		\item[until...until...] Yale has an important change here: ``until it is given unto him to know the mysteries of God, until \textit{they} know them in full.'' So the double ``until'' could be referring back to verse 9 with two different subjects: the first ``until'' refers to a single individual who gains knowledge of the mysteries, while the second ``until'' with ``they'' as subject refers to how the first individual will impart until the people he's sharing with also know the mysteries in full.
		Note that verse 11 also changes from plural to singular at the start in Yale.
		\end{compactdesc}
		
		\item[11$\vert$]
		\begin{compactdesc}
		\item[they that...to them...] Yale has, ``he that...to him...'', but it also retains the plural in ``until they know nothing...''
		\end{compactdesc}
		
		\end{compactdesc}

	\subsubsection{Alma 13} % *AL13

		\begin{compactdesc}

		\item[3$\vert$]
		\begin{compactdesc}
		\item[preparatory redemption] What is it? The new birth. See my notes in the interesting stories of the BOM section, on Alma 18--19, Ammon's experience with Lamoni. 
		\end{compactdesc}

		\end{compactdesc}

	\subsubsection{Alma 17} % *AL17

		\begin{description}

		\item[10$\vert$]
		\begin{compactdesc}
		\item[and they were comforted] They wanted to turn back; see Al 26:27.
		\end{compactdesc}
		
		\item[32$\vert$]
		\begin{compactdesc}
		\item[rushed forth with much swiftness] There's no record here of Ammon and his fellows receiving any help from the Lord in doing this---the first step they took was to exert themselves as much as they possibly could, on their own strength, apparently, in order to begin to fix the situation that other people had put them in. They had to sprint and really put forth all their effort in order to get the sheep back in the right spot.
		\end{compactdesc}

		\end{description}

	\subsubsection{Alma 18} % *AL18

		\begin{compactdesc}

		\item[15$\vert$]
		\begin{compactdesc}
		\item[they supposed that whatsoever they did was right] See Jud 17:6; 21:25;
		\end{compactdesc}
		
		\item[22$\vert$]
		\begin{compactdesc}
		\item[being wise, yet harmless] Cf. Christ talking to the apostles: ``Behold, I send you forth as sheep in the midst of wolves: be ye therefore wise as serpents, and harmless as doves'' (Mt 10:16).
		\end{compactdesc}

		\end{compactdesc}
		
	\subsubsection{Alma 19} % *AL19
	
		\begin{compactdesc}
		
		\item[30$\vert$]
		\begin{compactdesc}
		\item[she clasped her hands] See Yale, where it's ``clapped her hands.''
		\end{compactdesc}
		
		\end{compactdesc}

	\subsubsection{Alma 20} % *AL20

		\begin{compactdesc}

		\item[13$\vert$]
		\begin{compactdesc}
		\item[to his astonishment, his father was angry with him] This reminds me of Rodger Liddle's conversion story. When he told his wife that he had been reading the LDS scriptures and wanted to go to church, she became so angry that her face and look completely changed. He said it was like the possessed girl in the movie \textit{Exorcist}. He said he had never seen anyone filled with so much anger and hatred.
		\end{compactdesc}

		\end{compactdesc}

	\subsubsection{Alma 22} % *AL22

		Could this chapter be taking place inside the temple? They are in the palace of the land of Nephi (Al 22:1--2), which is where the temple was built (2N 5:7--8, 16). Aaron tells the king to ``bow down before God,'' which is what we do in the temple, before the altar. The king then bowed down ``before the Lord, upon his knees.'' Maybe this is taking place at an altar in the temple? Note that they talk about temples, many of them, in 23:2.

		\begin{compactdesc}

		\item[14$\vert$]
		\begin{compactdesc}
		\item[since man had fallen...] Can we read ``since'' as temporal here, and not causal? The natural reading is causal here, but it looks like most of the occurrences of ``since'' are temporal in the BOM. This reading makes sense though---before the fall, man was innocent and pure (DC 93:38), and since that time, he hasn't merited anything of himself. So can you merit something of yourself when you're pure and innocent? Perhaps you can. Does every rebirth include becoming pure and innocent again? Does that happen when we die in this life and move onto the next?
		\end{compactdesc}

		\item[18$\vert$]
		\begin{compactdesc}
		\item[as if he were dead] Again, the language death and dying when speaking of the new birth.
		\end{compactdesc}

		\item[20$\vert$]
		\begin{compactdesc}
		\item[mightier than us all] See Ab 3:19. 
		\end{compactdesc}

		\item[22$\vert$]
		\begin{compactdesc}
		\item[feared lest that...there should be a great contention] He probably had heard from Ammon and perhaps Abish what had happened when she rounded up the people, so he didn't want that to happen again. Then it starts to happen anyway in 24.
		\end{compactdesc}

		\item[27--35$\vert$]
		\begin{compactdesc}
		\item[Geographic digression] The king sends a proclamation, and then we get kind of a lengthy digression on Book of Mormon geography.
		\end{compactdesc}

		\end{compactdesc}

	\subsubsection{Alma 23} % *AL23

		\begin{description}

		\item[4$\vert$]
		\begin{compactdesc}
		\item[consecrating priests and teachers...among the Lamanites] They didn't seem to have a problem giving the priesthood to Lamanites, who supposedly were a ``cursed people'' and one separate from the Lord. As soon as they repented, they consecrated priests, gave them the priesthood, and started the church. 
		\end{compactdesc}

		\end{description}

	\subsubsection{Alma 24} % *AL24

		\begin{description}
		
		\item[3$\vert$]
		\begin{compactdesc}
		\item[called his name Anti-Nephi-Lehi] This has got to be a new name, right? He's got to be pretty old already, so the king is giving him a new name---he's taking the name of the whole people on himself.
			Also, why doesn't he give the kingdom to Lamoni?
		\end{compactdesc}

		\item[6$\vert$]
		\begin{compactdesc}
		\item[not one soul..would take up arms] That this seems strange to us is probably at least in part to cultural bias. Also, remember from 18:5--6 and 19:20 that killing was a hot topic among the people, and there were divided opinions about whether it was permissible. It seems like the people of Anti-Nephi-Lehi wanted to separate themselves completely from that practice and that discussion, so they decided not to kill at all.
		\end{compactdesc}
		
		\item[10$\vert$]
		\begin{compactdesc}
		\item[taken away the guilt from our hearts, through the merits of his Son] Shows (or at least seems to show) that there is an emotional aspect to repentance that doesn't have to do with wiping away responsibility for or accountability for sin. Repentance does that, but it doesn't necessarily follow that a person feels guilt about what they've done. If they do feel guilt, then the Atonement can help them with that too, even though strictly speaking that may not be a part of repentance.
		\end{compactdesc}

		\end{description}

	\subsubsection{Alma 26} % *AL26

		\begin{description}

		\item[22$\vert$]
		\begin{compactdesc}
		\item[reveal things which never have been revealed] This verse can't be talking about learning the doctrines of the gospel, because those have already been revealed. The verse is referring to things which have \textit{never} been revealed. What things are those? The things inside our hearts, and the things inside the hearts of other people.
			For more study of this verse, see it under the ``mysteries'' section.
		\end{compactdesc}

		\item[24$\vert$]
		\begin{compactdesc}
		\item[the ways of a transgressor from the beginning] Note that the Nephites fall into the same trap as the Lamanites, thinking that the Lamanites just are a certain way by their nature, and so they're hopeless to try to deal with. We see the Lamanites' view at Mo 10:12--13 and Al 20:10, 13.
		\end{compactdesc}

		\end{description}

	\subsubsection{Alma 29} % *AL29

	This is Alma's song. There are other ones in the scriptures, like of Nephi and Moses. This one comes from Alma. It is an unusual chapter---it seems out of place, the speaker isn't identified by name (11--15 suggest Alma; 9 and 16 sound more like Ammon), and it isn't given any commentary on either side. But it's very beautiful and has some important things in it.
	
		\begin{compactdesc}
		
		\item[1$\vert$]
		\begin{compactdesc}
		\item[trump of God] See ``trump'' in the concepts.
		\end{compactdesc}

		\item[3--5$\vert$]
		\begin{compactdesc}
		\item[Differences in Yale] The Yale edition changes the punctuation of these verses and the changes open up the way for new interpretations. Here are a few:
			\begin{itemize}
			\item ``I am a man and do sin in my wish'' (no comma)
			\item ``decrees which are unalterable according to their wills'' (no comma. The way I read 1981 is, the decrees are unalterable and are decreed because of the will of men. But there's a different reading, more natural in the Yale, which says that the decrees are unalterable ``according to'' or because of the will of men)
			\item ``...before all men---or he that knoweth not good from evil is blameless---but he that knoweth good and evil, to him...''
			\end{itemize}
		The moral of the story is to always check the Yale when you have a question. As for the exact meaning of this passage, I have no idea. I know what it is saying, but I also think it's saying something else that I am missing.
		\end{compactdesc}
		
		\item[5$\vert$]
		\begin{compactdesc}
		\item[to him it is given according to his desires] See Ps 106:13--15, especially 15.
		\end{compactdesc}

		\end{compactdesc}

	\subsubsection{Alma 30} % *AL30

	It's interesting to me that Korihor isn't mentioned in connection with Nehor. It seems like he would have been involved with that order.
	
		\begin{description}

		\item[29$\vert$]
		\begin{compactdesc}
		\item[they delivered him up] What authority do they have to do this? I don't think we're told that he's broken any laws. On what grounds do they hold him and send him to Zarahemla?
		\end{compactdesc}
		
		\item[44$\vert$]
		\begin{compactdesc}
		\item[denote] The only occurrence of this word in all the scriptures. There are two definitions in the 1828:
			\begin{enumerate}
			\item To mark; to signify by a visible sign; to indicate; to express.
			\item To show; to betoken; to indicate.
			\end{enumerate}
		Alma says, ``all things denote there is a God.'' He has also previously said that ``all things are a testimony that'' God exists and Christ will come (41). This verse, 44, deserves some careful attention, more than I can give it right now.
		\end{compactdesc}
		
		\begin{description}

		\item[59$\vert$]
		\begin{compactdesc}
		\item[he was run upon...until he was dead] There are other examples, especially in the OT, of people ``running upon'' other people to kill them---it seems to be just charging with a group of people, maybe with weapons. But why do the Zoramites do it? Why those people, in particular? What would they have had against Korihor, or in themselves, that would lead them to kill him?
		\end{compactdesc}

		\end{description}

		\end{description}
		
	\subsubsection{Alma 31} % *AL31

		\begin{compactdesc}

		\item[35$\vert$]
		\begin{compactdesc}
		\item[many of them are our brethren] Who else was there? A clue for more people being around? Or mixing with the Lamanites? Why did he include this remark? Maybe they were brothers in the church, so he's talking about them belonging to the church?
		\end{compactdesc}

		\end{compactdesc}

	\subsubsection{Alma 32} % *AL32

		\begin{description}
		
		\item[26--28$\vert$]
		\begin{compactdesc}
		\item[words/words] Plural ``words'' in 26--27, singular ``word'' in 28. Important shift; we're talking about different ``words'' here.
		\end{compactdesc}
		
		\item[27$\vert$]
		\begin{compactdesc}
		\item[arouse your faculties] A few quotes:
			\begin{compactdesc}
			\item[TPJS 14] Are not these things a sufficient stimulant to arouse the faculties, and call forth the energies of every man, woman or child, that possesses feelings of sympathy for their fellows, or that is in any degree endeared to the budding cause of our glorious Lord?
			\end{compactdesc}
		\end{compactdesc}
		
		\item[28--43$\vert$]
		\begin{compactdesc}
		\item[word/seed/tree] These words are interchangeable in these verses, I think.
		\end{compactdesc}

		\item[35$\vert$]
		\begin{compactdesc}
		\item[because it is light] See Eph 5:13, ``But all things that are reproved are made manifest by the light: for whatsoever doth make manifest is light.''
		\end{compactdesc}
		
		\item[42$\vert$]
		\begin{compactdesc}
		\item[tree] See Ps 1:1--3: ``BLESSED is the man that walketh not in the counsel of the ungodly, nor standeth in the way of sinners, nor sitteth in the seat of the scornful. 2 But his delight is in the law of the LORD; and in his law doth he meditate day and night. 3 And he shall be like a tree planted by the rivers of water, that bringeth forth his fruit in his season; his leaf also shall not wither; and whatsoever he doeth shall prosper.''
		\end{compactdesc}

		\end{description}
		
	\subsubsection{Alma 34} % *AL34
	
		\begin{description}
		
		\item[16$\vert$]
		\begin{compactdesc}
		\item[thus mercy can satisfy the demands of justice] In the individual case, I think. For individual people, not generally, here. 
		\end{compactdesc}
		
		\item[27$\vert$]
		\begin{compactdesc}
		\item[drawn out in prayer] See Hel 13:22.
		\end{compactdesc}
		
		\end{description}
		
	\subsubsection{Alma 37} % *AL37
	
		\begin{compactdesc}
		
		\item[6$\vert$]
		\begin{compactdesc}
		\item[small means...doth confound the wise] Eth 3:5, ``Behold, O Lord, thou canst do this. We know that thou art able to show forth great power, which looks small unto the understanding of men.''
		\end{compactdesc}
		
		\end{compactdesc}
		
	\subsubsection{Alma 38} % *AL38
	
		Why does Shiblon, who seems like the most obedient and diligent son, get the shortest address? It could be that he already heard a fuller version of his father's conversion on his missionary journeys; remember that Shiblon and Corianton went on the missionary trip in Al 31--35, but Helaman didn't.
		
	\subsubsection{Alma 41} % *AL41
		
		\begin{compactdesc}
		
		\item[7$\vert$]
		\begin{compactdesc}
		\item[stand or fall] See Rom 14:4, ``Who art thou that judgest another man's servant? to his own master he standeth or falleth. Yea, he shall be holden up: for God is able to make him stand.'' For the middle part, it's \textgreek{τῷ ἰδίῳ κυρίῳ στήκει ἢ πίπτει}.
		\end{compactdesc}
		
		\end{compactdesc}

\subsection{Helaman} % *HEL

	\subsubsection{Helaman 5} % *HEL5
	
		\begin{compactdesc}
		
		\item[12$\vert$]
		\begin{compactdesc}
		\item[shafts in the whirlwind] See 3N 8:16.
		\end{compactdesc}
		
		\item[21--49$\vert$]
		\begin{compactdesc}
		\item[The story] Here's one way to describe this spiritual manifestation, and it works for some other manifestastions in the scriptures as well: there are times when our spiritual standing, and the spiritual standing of others, is made manifest to our regular senses, not just to a spiritual sense. Here we see the spiritual standing of Nephi and Lehi: pure and clean, refined and cleansed by fire. And we see the spiritual standing of those who oppress them: ``they were overshadowed with a cloud of darkness, and an awful solemn fear came upon them.'' So sometimes the Spirit quickens our faculties enough for us to perceive holiness and wickedness with the regular old five senses.f
		\end{compactdesc}
		
		\end{compactdesc}
		
	\subsubsection{Helaman 6} % *HEL6
	
		\begin{compactdesc}
		
		\item[21--31$\vert$]
		\begin{compactdesc}
		\item[The robbers' covenants] There is a lot of language here that is very similar to priesthood language. Look especially at 21--24; you don't have to change the description very much to make this sound like the language of the temple and the priesthood covenants made there. My guess is that lots of people have connected this passage to Masonry, which (from what I understand) went through a period of disfavor maybe around the 1810's and 1820's, but began to grow in popularity again after that, until eventually Joseph joined in the 1840's (I think he joined very soon before restoring the endowment). So there may be hints at Masonry here, but the point isn't restricted to that group, since there's more to be said here about corruption in general.
		\end{compactdesc}
		
		\end{compactdesc} 
		
	\subsubsection{Helaman 8} % *HEL8
	
		The point of this chapter is to prove to the people that there are prophets, that they prophesy of Christ, and that Nephi is one of those. Nephi cites a bunch of examples from the scriptures of both prophetes and true prophecies, and then points out that they people can't really deny these examples. He asks, so if they can prophesy of those things, why not of the coming of Christ? Then he connects their wickedness to the death of the chief judge, and he makes his own prophecy about that. So he closes this short discourse with a prophecy of his own to prove that he himself is a prophet. Very nice chapter and well integrated I think.
		
		The people listening to him clearly understand his words in this way, as we see in 9:1--2, where men checking the chief judge see what they are doing as a test of Nephi's prophethood. Also 9:16, and 23.
		
	\subsubsection{Helaman 9} % *HEL9
	
		The second half of this chapter is another test of Nephi's prophethood, just like the previous chapter.
	
		\begin{compactdesc}
		
		\item[5$\vert$]
		\begin{compactdesc}
		\item[when they saw they believed] This is a very powerful moment. It's the moment of realizing something new, understanding something, being \textit{compelled} to believe something for the first time. These people cannot deny now what they have seen; they can't erase this experience. They are forced to recognize that Nephi is a prophet and that the other things he is saying are true as well. 
		
		We have these moments too, when we witness something or have some experience which leads to a belief that we cannot deny. This is a very direct example of a sensory witness, or a visual experience that---in connection with their other recent experiences with Nephi---they cannot deny. It doesn't have to be visual though. Spiritual realizations and connections can happen in the very same way and be just as powerful, totally undeniable and incontrovertible by those who experience them.
		
		See 1N 4:16--17 for a similar moment of clarity, that one much more spiritual.
		\end{compactdesc}
		
		\item[41$\vert$]
		\begin{compactdesc}
		\item[[he] has told us things] Very interesting use of ``things'' here. It has to mean something like, ``events in the world.'' See also Hel 15:8; % *THINGS
		\end{compactdesc} 
		
		\end{compactdesc} 
		
	\subsubsection{Helaman 10} % *HEL10
	
		\begin{compactdesc}
		
		\item[1--11$\vert$]
		\begin{compactdesc}
		\item[Receiving power from the Lord] I'm not sure how to describe the powers Nephi receives. He gets one kind of power in verse 5, or maybe various kinds, because he's going to be might in word and in deed, in faith and in works. All things will be done according to his word (the repetition of ``word'' here is probably significant). So this seems to be one power that he gets, similar to the one mentioned in Jac 4:6, about commanding the elements of the earth and them obeying.
		
		Verse 6 talks again about this same power, and then in 7 the Lord actually bestows the power. He describes it as a power of sealing and loosing on earth and in heaven, which makes it seem like what we would call the ``sealing'' power, but then the Lord says at the end of the verse, ``and thus shall ye have power among this people.'' The \textit{thus} is the important part, and so is the \textit{thus} at the beginning of verse 8. I think that those words make it clear that the Lord has in mind only a single power here, and that whatever the ``sealing'' power is, it includes power over people and the earth in the way the Lord describes here.
		
		Also, another interesting thing---the Lord promises the power in verses 5 and 6, then actually gives the power in verse 7. So even though these two things happen very close together in time, they are still separate---the promising of the power comes first, and then the giving of the power.
		\end{compactdesc}
		
		\end{compactdesc}
		
	\subsubsection{Helaman 11} % *HEL11
	
		\begin{compactdesc}
		
		\item[19$\vert$]
		\begin{compactdesc}
		\item[not a whit behind him] So I think we can assume that Lehi had received the same powers that Nephi had, maybe?
		\end{compactdesc}
		
		\end{compactdesc}
		
	\subsubsection{Helaman 16} % *HEL16
	
		\begin{compactdesc}
		
		\item[18$\vert$]
		\begin{compactdesc}
		\item[not reasonable] I love this phrase---``it doesn't make sense,'' ``it isn't possible,'' etc. They think themselves out of it. I guess what puzzles me about this way of thinking is that it seems to grant us a very great deal of knowledge and a very powerful epistemic perspective---it's the idea that we, from where we are right now, are capable of judging whether such a being as a Christ is reasonable/possible or not. What reason do we have to think that we are in such a position, that we're that capable of making that judgment? Not sure.
		\end{compactdesc}
		
		\item[21$\vert$]
		\begin{compactdesc}
		\item[we depend upon them to teach us the word] I think this is one of Joseph Smith's greatest innovations---teaching that each person has the power to receive the revelations and the assurances and all the knowledge. He fixed this problem by teaching that everyone can have it all.
		\end{compactdesc}
		
		\end{compactdesc}
		
\subsection{3 Nephi} % *3N

	\subsubsection{3 Nephi 1} % *3N1
	
		\begin{compactdesc}
		
		\item[9$\vert$]
		\begin{compactdesc}
		\item[should be put to death] What authority did they have to enforce this? It's kind of strange. The government must have been very corrupt to allow something like this. See also 6:22--25, which makes this even harder to understand. It's odd that there's no comment about it.
		\end{compactdesc}
		
		\item[14$\vert$]
		\begin{compactdesc}
		\item[I come unto my own] His own what? People? Creations? Kinds of beings (humans)?
		\end{compactdesc}
		
		\item[22$\vert$]
		\begin{compactdesc}
		\item[that they might not believe] It would be hard to explain away a missing night, I think. 
		\end{compactdesc}
		
		\item[24--25$\vert$]
		\begin{compactdesc}
		\item[Dispute over the law of Moses] It says that ``the word came unto them'' in order to resolve the disagreement and fix the misunderstanding. So it almost sounds like a revelation was necessary to fix this problem. 
		\end{compactdesc}
		
		\item[28$\vert$]
		\begin{compactdesc}
		\item[many dissenters] Why leave at this point? Doesn't seem to make sense. Why would they jump ship after the clear and obvious signs? Maybe they wanted money or something.
		\end{compactdesc}
		
		\end{compactdesc}
		
	\subsubsection{3 Nephi 2} % *3N2

		\begin{compactdesc}
		
		\item[1$\vert$]
		\begin{compactdesc}
		\item[less and less astonished] I guess the idea is that signs and wonders don't feed faith over the long term. Only receiving the Holy Ghost will do that.
		\end{compactdesc}
		
		\item[1--2$\vert$]
		\begin{compactdesc}
		\item[Disbelieving] I think this is what happens to some people who leave the church. They have to explain away the spiritual experiences they've had; they have to make those experiences into nothing or else their decision to abandon their testimony will look foolish. I like that the text here says that they imagine a ``vain thing'' in their hearts.
		\end{compactdesc}
		
		\item[8$\vert$]
		\begin{compactdesc}
		\item[began to reckon their time...] It's strange that they changed their way of counting time but no longer believed the events that marked the change. It seems like the changed way of counting time would be a constant reminder of what had happened.
		\end{compactdesc}
		
		\item[12$\vert$]
		\begin{compactdesc}
		\item[their freedom and their liberty] There may be something to this difference in the BOM, between freedom and liberty.
		\end{compactdesc}
		
		\end{compactdesc}
		
	\subsubsection{3 Nephi 3} % *3N3
	
		\begin{compactdesc}
		
		\item[4$\vert$]
		\begin{compactdesc}
		\item[unconquerable spirit] Al 52:33, ``And it came to pass that Jacob, being their leader, being also a Zoramite, and having an unconquerable spirit, he led the Lamanites forth to battle with exceeding fury against Moroni.''
		\item[everlasting hatred] One other time in the scriptures, Al 37:32: ``And now, my son, remember the words which I have spoken unto you; trust not those secret plans unto this people, but teach them an everlasting hatred against sin and iniquity.''
		\item[many wrongs] Giddianhi seems to be taking up the banner of the past Lamanite sufferings, as previous Lamanites have done, but that's kind of strange because there are a lot of Nephites among them (3N 1:28). He explains himself in verse 10 though.
		\end{compactdesc}
		
		\end{compactdesc}
		
	\subsubsection{3 Nephi 5} % *3N5
	
		\begin{compactdesc}
		
		\item[8$\vert$]
		\begin{compactdesc}
		\item[in the eyes of some] Interesting comment here. Why would only some find these things great and marvelous? Because they actually aren't? 
		\end{compactdesc}
		
		\item[12$\vert$]
		\begin{compactdesc}
		\item[after their transgression] What is this transgression? The first time the people of Nephi (or a group from it) turned away from the Lord as a whole?
		\end{compactdesc}
		
		\item[14$\vert$]
		\begin{compactdesc}
		\item[who were the holy ones] So who are we talking about here? Prophets like Nephi, Benjamin, Alma, I'm guessing. I wonder why he points out that they were holy, though. Maybe to put emphasis on the fact that they had \textit{become} holy, or that they were holy and so their prayers were more likely to be answered.
		\end{compactdesc}
		
		\end{compactdesc}
		
	\subsubsection{3 Nephi 6} % *3N6
	
		\begin{compactdesc}
		
		\item[3$\vert$]
		\begin{compactdesc}
		\item[who were desirous to remain Lamanites] It's hard to get a handle on what distinguishes the Lamanites and Nephites sometime. The groups move back and forth, and some seem to become part of the other, but here the distinction is made clear again and it has to do with a personal decision made by some of the people.
		\end{compactdesc}
		
		\end{compactdesc}
		
	\subsubsection{3 Nephi 9} % *3N9
	
		\begin{compactdesc}
		
		\item[17$\vert$]
		\begin{compactdesc}
		\item[given to become] This construction occurs several other times in the scriptures, like at Jn 1:12, DC 11:30, and something similar in DC 45:8. But in those places you get the word ``power''---we receive \textit{power} to become the sons of God. Why omit the word here? What's the significance of leaving it out?
		\end{compactdesc}
		
		\item[20$\vert$]
		\begin{compactdesc}
		\item[with fire and with the Holy Ghost] Are these the same? Or different? Why ``and'' connecting them? I think I've thought about this before and have had an answer, but I can't think of it right now. See DC 19:31.
		\end{compactdesc}
		
		\end{compactdesc}
		
	\subsubsection{3 Nephi 11} % *3N11
	
		When do these events happen? That is, how long after the destruction and prophecy give in chapters 8--10? 8:5 says that the destruction started on the fourth day of the first month of year 34. Then in 10:18 Mormon says that Christ manifested himself to the people ``in the ending of the thirty and fourth year.'' So nearly a year? At least nine months? Some decent amount of time; it wouldn't have been right after, it seems. Note that it's definitely ``after the ascension'' into heaven (10:18), which is recorded in Acts 1.
		
		\begin{compactdesc}
		
		\item[1$\vert$]
		\begin{compactdesc}
		\item[great and marvelous change] I've heard a few interpretations for this phrase: the change in the land is one (due to the destruction), and the change in their hearts is another. I think the second interpretation is very beautiful. They were marveling and wondering with each other about how different things were, and how differently they felt about everything. It could also be just how different their lives are now, with so many of their people dead and the cities changed and everything.
		\end{compactdesc}
		
		\item[8$\vert$]
		\begin{compactdesc}
		\item[they thought it was an angel] Kind of strange, since the voice just told them to behold the beloved son.
		\end{compactdesc}
		
		\item[14$\vert$]
		\begin{compactdesc}
		\item[thrust/feel] See Jn 20:25.
		\end{compactdesc}
		
		\item[21$\vert$]
		\begin{compactdesc}
		\item[power to baptize] Why did the Lord give Nephi power to baptize, when he had already been doing baptisms? (A bunch of verses show that, but see 3N 1:23, or 8:24--25, for example.) It could be that, with the turning over of things from the law of Moses to the law of Christ, and the gospel being given in a different way, that the Lord needed to give new authority for the new way of doing things (see 3N 12:46--47, especially 47). See verse 25, for example; compare to the baptism in Mo 18:13. Alma's baptism just seems to mark a covenant with God to serve him (see Mo 21:35 as well, and probably others), whereas this new baptism, instructed by Christ, names the godhead and names Christ himself as the giver of the authority. So it could be that there's something different about the baptism itself as well. That's all I can think of right now at least. For more about the older covenant, see, for example, 2Ch 15:12;
		\end{compactdesc}
		
		\item[32--33$\vert$]
		\begin{compactdesc}
		\item[repent, believe, be baptized] So everyone, everywhere, is command to repent and believe in Christ; but only those who are baptized will be saved. See the topic for the universality of the gospel---people can be commanded to repent and believe in Christ in many different ways, and their repentance and faith can take many different forms (think of those things as ordinances as well, faith and repentance). But not baptism, that's different. Baptism has to be done in the correct way. That could be why Christ is really emphasizing that there can't be disagreements over how baptism should be done.
		\end{compactdesc}
		
		\item[37$\vert$]
		\begin{compactdesc}
		\item[Shift in pronouns] Very important; he's been talking about ``whoso'' (whoever), and now speaks directly to the people. See my notes about this verses and the ones preceding it at ``universality of the gospel.''
		\end{compactdesc}
		
		\item[37--38$\vert$]
		\begin{compactdesc}
		\item[Differences] The only one is at the end: ``receive these things'' vs. ``inhereit the kingdom of God.'' Christ repeats himself for emphasis here.
		\end{compactdesc}
		
		\item[39$\vert$]
		\begin{compactdesc}
		\item[Shift in pronouns again] Back to ``whoso.''
		\end{compactdesc}
		
		\end{compactdesc}
		
	\subsubsection{3 Nephi 12} % *3N12
	
		\begin{compactdesc}
		
		\item[1--2$\vert$]
		\begin{compactdesc}
		\item[see and know/believe on words] There are people who see and know for themselves, and people who believe on the words of those first people; see DC 46:13--14, but especially Al 32:13--19, where there's a parallel with humbling yourseld and being obedient, and knowing the will of God vs. only having a reason to believe.
		\end{compactdesc}
		
		\item[23--24$\vert$]
		\begin{compactdesc}
		\item[Differences in Matthew] On the idea of faith and repentance being ordinances---note the changes in these two verses between Matthew and 3 Nephi. The language about bringing a gift to the altar---very easy to interpret that as an action of an ordinance---is changed to something about coming to Christ, or desiring to do so. So the clear ordinance language is changed to language that isn't so clear for ordinances, but I think it's a hint that we should see it that way; see also 3N 9:20--22, and the earlier verse 19 of chapter 12. The gift we offer at the altar is a broken heart and a contrite spirit; ``coming unto Christ'' is believing in him and repenting; and we can see those actions as ordinances, with the offering of our heart and spirit as the action accompanying it. The connection to 3N 9:22 is important too, with the stuff about being like a child; that helps us explain what a broken heart and contrite spirit is.
		\end{compactdesc}
		
		\item[37$\vert$]
		\begin{compactdesc}
		\item[cometh of more] Interesting difference between this verse and the corresponding one in Matthew; I have no idea what it could mean right now. Whatever things \textit{arise} out of more complex communication are evil? That doesn't really make sense. Maybe just, whatever comes out of dishonest communication is evil; the yay yay nay nay stuff is just about simple, honest, and direct speech. Maybe.
		\end{compactdesc}
		
		\item[45$\vert$]
		\begin{compactdesc}
		\item[Difference with Matthew] Deletes the part about sending rain on the just and unjust. Interesting.
		\end{compactdesc}
		
		\item[48$\vert$]
		\begin{compactdesc}
		\item[Father who is in heaven] Could refer to Christ, so this could be an explanatory ``or''; this is a very good example of why looking at the word ``or'' would be really interesting in the scriptures.
		\end{compactdesc}
		
		\end{compactdesc}
		
	\subsubsection{3 Nephi 13} % *3N13
	
		\begin{compactdesc}
		
		\item[25--34$\vert$]
		\begin{compactdesc}
		\item[Who receives these remarks?] The BOM makes it clear that Christ is speaking directly to the disciples, the twelve disciples, in these verses. In 14:1 he returns to spekaing with the multitude.
		\end{compactdesc}
		
		\end{compactdesc}
		
	\subsubsection{3 Nephi 14} % *3N14
	
		\begin{compactdesc}
		
		\item[1$\vert$]
		\begin{compactdesc}
		\item[Who receives these remarks?] The BOM makes it clear that Christ now returns to addressing the entire multitude here.
		\end{compactdesc}
		
		\end{compactdesc}
		
	\subsubsection{3 Nephi 17} % *3N17
	
		\begin{compactdesc}
		
		\item[16--17$\vert$]
		\begin{compactdesc}
		\item[saw Jesus speak] 16 says, ``the eye hath never seen...so great and marvelous things as we saw...Jesus speak unto the Father.'' What did they \textit{see} him speak to the Father? They \textit{heard} his words; but what did they \textit{see}? Note 3N 18:24.
		\end{compactdesc}
		
		\end{compactdesc}
		
	\subsubsection{3 Nephi 18} % *3N18
	
		\begin{compactdesc}
		
		\item[28$\vert$]
		\begin{compactdesc}
		\item[knowingly] Who's the one doing something knowingly here---the people partaking of the sacrament, or those who are letting it happen? See next verse---those who allow.
		\end{compactdesc}
		
		\item[30$\vert$]
		\begin{compactdesc}
		\item[repents and is baptized] This verse makes it seem like the ``unworthy'' partaking from the previous two verses is by someone outside of the church. But people inside, already baptized, must be able to partake unworthily as well.
		\end{compactdesc}
		
		\item[38--39$\vert$]
		\begin{compactdesc}
		\item[cloud overshadowed them] After everything the people had just seen and experienced, why did a cloud descend to block them from seeing this? It seems that the disciples saw it, but not the larger group.
		\end{compactdesc}
		
		\end{compactdesc}
		
	\subsubsection{3 Nephi 19} % *3N19
	
		Christ offers three prayers in this chapter. What is the significance of that number? And why is it the third prayer only that we don't get a record of (that we \textit{can't} have a record of?)
		
		Note also that as the people continue to pray under the influence of the Spirit, their spiritual experiences continue to get better and better.
	
		\begin{compactdesc}
		
		\item[25$\vert$]
		\begin{compactdesc}
		\item[transfiguration] It's something like the veil being taken from people so that they can see the proper glory of Jesus; they can perceive correctly for a moment, and they are purified by the glory and power emanating from the Lord. Reminds me of DC 76:118, ``that through the power and manifestation of the Spirit, while in the flesh, they may be able to bear his presence in the world of glory.''
		\item[countenance did smile upon them] Think of the smile as the ``mechanism'' of the transfiguration or purification? What happens when Christ smiles? He's happy, he's pleased, he sees people doing things he wants them to do and that he can't \textit{force} them to do; so his desires are filled, he becomes happy, and this light spills out of his countenance. This is a crude idea but there's something important going with his smiling. See also verse 30.
		\item[upon earth] Nothing \textit{on earth} could be this white; so being transfigured is like being off earth for a moment.
		\end{compactdesc}
		
		\end{compactdesc}
		
	\subsubsection{3 Nephi 20} % *3N20
	
		\begin{compactdesc}
	
		\item[24$\vert$]
		\begin{compactdesc}
		\item[Samuel] Samuel the \textit{Lamanite}, that is. At first I thought it meant Samuel in the OT.
		\end{compactdesc}
		
		\end{compactdesc}
		
	\subsubsection{3 Nephi 27} % *3N27
	
		\begin{compactdesc}
		
		\item[24$\vert$]
		\begin{compactdesc}
		\item[...of that which hath been] This is an unusual verse; not quite sure what to make of it right now.
		\end{compactdesc}
		
		\end{compactdesc}
		
	\subsubsection{3 Nephi 28} % *3N28
	
		\begin{compactdesc}
		
		\item[2$\vert$]
		\begin{compactdesc}
		\item[that we may speedily come unto thee] In other words, skip the waiting before the resurrection, and join Christ in his kingdom immediately with a body (I think). But note that he says more to the three other disciples about the change in the ``twinkling'' of an eye.
		\end{compactdesc}
		
		\item[7--9$\vert$]
		\begin{compactdesc}
		\item[Changes in the body to live longer] Three changes happen here:
			\begin{compactitem}
			\item Not taste of death and live a very long time.
			\item Not endure the pains of death and be changed to immortality instantly.
			\item Not have pain or sorrow in the flesh.
			\end{compactitem}
				See also:
					DC 7:1--3, 6: ``AND the Lord said unto me: John, my beloved, what desirest thou? For if you shall ask what you will, it shall be granted unto you. 2 And I said unto him: Lord, give unto me power over death, that I may live and bring souls unto thee. 3 And the Lord said unto me: Verily, verily, I say unto thee, because thou desirest this thou shalt tarry until I come in my glory, and shalt prophesy before nations, kindreds, tongues and people...6 Yea, he has undertaken a greater work; therefore I will make him as flaming fire and a ministering angel; he shall minister for those who shall be heirs of salvation who dwell on the earth.''
		\end{compactdesc}
		
		\item[9$\vert$]
		\begin{compactdesc}
		\item[shall not have pain] So there is some kind of change that can be effected in our physical bodies such that we won't experience pain. Pain is normally an indicator of bodily harm, though, so it's important that the disciples have already been promised that they won't die. Since they don't die, they don't need to feel pain. 
		\end{compactdesc} 
		
		\end{compactdesc}
		
\subsection{Mormon} % *MORMON

	\subsubsection{Mormon 1} % *MORM1
	
		\begin{compactdesc}
		
		\item[2$\vert$]
		\begin{compactdesc}
		\item[quick to observe] I love David Bednar's talk about this topic.
		\end{compactdesc}
		
		\end{compactdesc}
		
	\subsubsection{Mormon 2} % *MORM2
	
		\begin{compactdesc}
		
		\item[1$\vert$]
		\begin{compactdesc}
		\item[the people of Nephi appointed me] Why did they appoint such a young kid? They could see how righteous he was; righteous action makes you seem very wise. They knew that he had a spiritual power that they didn't have. Maybe they wouldn't have identified it like that, but they knew there was something different about him.
		\end{compactdesc}
		
		\item[14$\vert$]
		\begin{compactdesc}
		\item[curse God, and wish to die] See Job 1:5; 2:9;
		\end{compactdesc}
		
		\end{compactdesc}
		
	\subsubsection{Mormon 5} % *MORM5
	
		\begin{compactdesc}
		
		\item[16$\vert$]
		\begin{compactdesc}
		\item[in the world] Related to the exchange I had with President Openshaw about the sons of perdition, it may be that the people spoken of here in Mormon have completely driven the Spirit and the light of Christ away from their lives \textit{in this world}, but out of this world, they would still have some of it. Until they return to their native element, that is?
		\end{compactdesc}
		
		\end{compactdesc}
		
	\subsubsection{Mormon 9} % *MORM9
	
		\begin{compactdesc}
		
		\item[23--24$\vert$]
		\begin{compactdesc}
		\item[Quoting Mk 16:15--18] What's interesting is that these verses in Mark are usually considered spurious and are taken out of the critical edition.
		\end{compactdesc}
		
		\end{compactdesc}
		
\subsection{Ether} % *ETHER

	\subsubsection{Ether 2} % *ETH2
	
		\begin{compactdesc}
		
		\item[9$\vert$]
		\begin{compactdesc}
		\item[ripened in iniquity] See also the next verse, ``fulness of iniquity.'' My question is, why use the word ``iniquity'' and not ``sin''? Not sure why. I haven't looked at the difference between them in a long time. See verse 15 though, where sin does make an appearance.
		
		I am impressed by these two verses (9--10) though that they give emphasis to the idea that it's the fulness of iniquity which will ultimately be punished; whatever the fulness is, that's what's required. See 11 and 15 also.
		\end{compactdesc}
		
		\item[19$\vert$]
		\begin{compactdesc}
		\item[no light; whither shall we steer?] This seems like a weird thing to ask---there's no light, so you ask where you're going to steer? They're \textit{inside} the boat; how are they going to know how to steer anyway, whether or not they have light inside with them or not? The boat is enclosed.
		
		So in this these barges might have been different from the ones they built previously (6, 16).
		\end{compactdesc}
		
		\end{compactdesc}
		
	\subsubsection{Ether 3} % *ETH3
	
		\begin{compactdesc}
		
		\item[1$\vert$]
		\begin{compactdesc}
		\item[carry them in his hands] President Openshaw used to talk about how this meant he was showing them great care, being very careful. It also means that they were pretty small, if he was carrying all 16 in his hands.
		\end{compactdesc}
		
		\item[13$\vert$]
		\begin{compactdesc}
		\item[because thou knowest these things] What are ``these things''? It could be the thing the boj has just said, about how God doens't lie, but it could also be whatever he learns in the moment when Christ reveals himself to him, because the Lord does that and \textit{then} says ``because thou knowest these things...'' See also 19, ``because of the knowledge of this man he could not be kept from beholding within the veil...'' Then see 26 as well, ``he knew that the Lord could show him all things.'' So the boj knew various things here, and the antecedent of the ``these things'' isn't just one thing, but a whole bunch of them.
		\item[the fall] Note that the boj says something about the fall himself in verse 2.
		\end{compactdesc}
		
		\item[15$\vert$]
		\begin{compactdesc}
		\item[never have I showed myself unto man whom I have created] This seems completely false; if you play around with the word ``man,'' it might make sense, as if you thought that the boj wasn't really a ``man'' any longer because he's been brought back into the presence of the Lord. Possible but weird. You might have to do that for other scriptures as well, including some mentioned in the Bible and maybe the DC, since if Christ revealed himself to Adam, then this is still going to be false. Weird comment.
		And then you get to verse 18, where Moroni reiterates that the boj is a ``man,'' and it gets weirder.
		\end{compactdesc}
		
		\end{compactdesc}
		
	\subsubsection{Ether 4} % *ETH4
	
		\begin{compactdesc}
		
		\item[7$\vert$]
		\begin{compactdesc}
		\item[revelations unfolding in later days] I mean ``later,'' not ``latter.'' When I read a verse like this, I think of how amazing sources have come to light for many people in these days, like Joseph Smith Papers project or CDBY or the writings of Eliza Snow or a whole bunch of other things. This has to be a partial fulfillment at least of these prophecies, so I see these new sources as windows into the kinds of visions that the boj had when Christ manifested himself to him. We get to peek at least a bit into what he saw and what he might have been thinking when we read JS on King Follett, for example, or BY on Adam-God.
		\end{compactdesc}
		
		\end{compactdesc}
		
	\subsubsection{Ether 5} % *ETH5
	
		\begin{compactdesc}
		
		\item[6$\vert$]
		\begin{compactdesc}
		\item[authority] What is Moroni talking about here? What authority? Speaking by the power and authority of the spirit? Or does he have keys over some part of the administration of God's work, something to do with the scriptures? Well we know that he has that, since he's the one who deals with Joseph Smith and the plates. So I think he does have some kind of authority over the book of Mormon or those records.
		\end{compactdesc}
		
		\end{compactdesc}

	\subsubsection{Ether 10} % *ETH10

		\begin{compactdesc}
		
		\item[11$\vert$]
		\begin{compactdesc}
		\item[but not unto himself] See 2Ch 25:2, ``And he did that which was right in the sight of the LORD, but not with a perfect heart.''
		\end{compactdesc}
		
		\end{compactdesc}
		
	\subsubsection{Ether 12} % *ETH12
	
		\begin{compactdesc}
		
		\item[12$\vert$]
		\begin{compactdesc}
		\item[my grace is sufficient for all men that humble themselves before me] 
		\end{compactdesc}
		
		\end{compactdesc}
		
\subsection{Moroni} % *MOR
	
	\subsubsection{Moroni 7} % *MOR7
	
		\begin{compactdesc}
		
		\item[19$\vert$]
		\begin{compactdesc}
		\item[lay hold] From the 1828 dictionary: ``It is much used after the verbs to take, and to lay; to take hold or to lay hold is to seize. It is used in a literal sense; as to take hold with the hands, with the arms, or with the teeth; or in a figurative sense.''
		\end{compactdesc}
		
		\item[24$\vert$] % *MOR7:24
		\begin{compactdesc}
		\item[there were divers ways...] How have I never noticed this before? This verse is incredible. Goes with levels principle, goodness, universality. Beautiful. 
			Note the change in Yale, removing the comma: ``there were divers ways that he did manifest things unto the children of men which were good.'' The ``which'' possibly has an unclear antecedent, between ``things'' and ``children of men,'' but given what's going on in this passage, the most natural reading is ``things.''
		\end{compactdesc}
		
		\item[42--44$\vert$]
		\begin{compactdesc}
		\item[must needs] Take 42 for example: ``if a man have faith he must needs have hope.'' I don't think that we are talking about conceptual entailments here. If we were talking about conceptual entailments, that would mean that, given how the concept of faith is, it would be necessary that we also have hope if we have faith. I don't think that's what Mormon is trying to get at here. Obviously the concepts of faith, hope, and charity are related to each other, but I don't think that we should understand these verses as talking about conceptual entailments among them. I would see this ``must needs'' language as an injunction to seek out faith, hope, and charity, or the one of those that you are missing; that having one isn't going to do you any good, so you ``must needs'' have the others. 44, for example, says that unless one is meek and lowly of heart, ``his faith and hope is vain.'' So he may still have faith and hope, but those things will be empty, useless for his salvation. I think that's a better way to think of it, rather than trying to work out a system where there are conceptual entailments among faith, hope, and charity such that, when we truly have one, we automatically have the others. It seems obvious you can have one without the others, depending on how you set them up.
		Same for the end of 44: someone can confess, even by the power of the Holy Ghost, that Jesus is the Christ, but unless he has charity, that confession means nothing to him.
		\end{compactdesc}
		
		\item[43--44$\vert$]
		\begin{compactdesc}
		\item[meek and lowly of/in heart] See Mor 8:25--26.
		\end{compactdesc}
		
		\item[45$\vert$]
		\begin{compactdesc}
		\item[believeth all things] Not realy sure what this means. Why would you believe \textit{all} things? What is the qualifier supposed to be? All things that one receives by revelation from God? All things from the prophet?
		\end{compactdesc}
		
		\item[46$\vert$]
		\begin{compactdesc}
		\item[charity never faileth] Could mean at least two things: when we have charity, we never fail others; because we have charity, God won't fail us.
		\item[greatest of all] Greatest of all what? Greatest of faith/hope/charity? I think it means, the greatest of all thing good things---recall Mor 7:28, ``they who have faith in him will cleave unto every good thing.'' Then we're told to cleave unto charity, ``which is the greatest of all [of the good things].''
		\item[for all things must fail] We can take it generally, as in it refers really to all things; that might be the most natural way. But if the ``greatest of all'' part previous to this refers to the greatest of all good things, then this part could mean that all good things except charity will fail. That's a little scary. But see the next verse, where we learn that charity endures forever---this ``fail'' thing must just mean that all other good things will pass away, not necessarily fail. Actually I like that interpretation. The first definition of ``fail'' in the 1828 is ``To become deficient; to be insufficient; to cease to be abundant for supply; or to be entirely wanting.'' It isn't until you get to definition \#10 that you get something more like how we normally use the term fail, to fail someone or let them down or disappoint them. So fail here just means that everything else you have will come up short in getting you to the presence of God---``for if he have not charity he is nothing'' (44), ``except men shall hav charity they cannot inherit that place which thou hast prepared in the mansions of thy Father'' (Eth 12:34).
		\end{compactdesc}
		
		\item[47$\vert$]
		\begin{compactdesc}
		\item[of Christ] Would be the ineluctable modality of sub/obj genitive, but see Eth 12:34: ``I know that this love which thou hast had for the children of men is charity.'' So you could still interpret this as objective genitive but subjective is probably better.
		\item[possessed of it] I think I've normally taken this to mean, ``whoever possesses charity,'' but you could take the ``of'' to mean ``by'' and think of it as, ``whoever is possessed by charity.'' This would be inline with Al 34:34--35, which talk about different spirits possessing our bodies at the time we go out of this life. Charity can possess us as well---we get the word in us, it infuses our bodies, and so we are possessed by that power. Note also that the two other occurrences of ``possessed of,'' both in the NT (Mt 8:33, Lk 8:36) are passive verbs.
		\end{compactdesc}
		
		\item[48$\vert$]
		\begin{compactdesc}
		\item[ye may be filled with this love] TPJS 9: ``Until we have perfect love we are liable to fall and when we have a testimony that our names are sealed in the Lamb's book of life we have perfect love and then it is impossible for false Christs to deceive us.''
		\item[we may be purified even as he is pure] See 3N 19:25 and 28, where the people are transfigured and become as white as Jesus's clothing. 28 notes that this is purification, and that it happens ``in Christ.''
		\end{compactdesc}
		
		\end{compactdesc}
		
	\subsubsection{Moroni 10} % *MOR10
		
		\begin{compactdesc}
		
		\item[3--5$\vert$]
		\begin{compactdesc}
		\item[Moroni's promise] Notes:
			The promise presupposes that you believe in God; otherwise how could you ponder on the mercy the Lord has shown to his children (verse 3), or how could you ask God (verse 4)?
				There is precedent in the BOM for asking when you don't understand very much---Lamoni and his father fit that description I think. But in any case the promise assumes a faith that some people may not have.
			The presupposition of faith is there again at the end of verse 4. After we're told what to ask, Moroni lists some qualifiers for our request: IF we ask with a sincere heart; IF we ask with real intent; IF we have faith in Christ.
				6--7 is interesting to read in light of this presupposition, for here we're told that we can know that Christ is by the power of the Holy Ghost, the same power that's going to manifest to us the truth of the BOM.
		\end{compactdesc}
		
		\end{compactdesc}

\section{Doctrine and Covenants} % 

	\subsubsection{DC 5} % *DC5 
	
		Many changes were made to this section from the time it was received, to the 1835 DC, and up until now. See JSR 28--31 for a summary and discussion. The most important change that I can see is that in the original revelation, Joseph was told that the gift to translate the BOM \textit{would be his only gift}, and that he would not receive another; the bits about ``for I will grant unto you no other gift until it is finished'' are not in the original. So this was revised later, because clearly Joseph either did receive other gifts or felt like he needed to---Marquardt in JSR suggests it was due to the Bible translation. In any case a lot of this language was changed to open up the possibility that he would receive other gifts.
		
		Other changes were made as well, some significant.
		
	\subsubsection{DC 10} % *DC10
	
		The date of this revelation is problematic; the given date in the current DC, summer of 1828, seems wrong. Should be around May 1829. See JSP-D1:38--39, and HDDC 200--205.
		
	\subsubsection{DC 18} % *DC18
	
		The date given is June 1829, and it's probably in the first half of the month, maybe the first 10 days or so. On June 14 of that year Oliver Cowdery wrote a letter to Hyrum Smith that includes a lot of language from this revelation, so it must have been received before that. The letter is quoted in HDDC 264, 267 and also appears at EMD 2:402--403.
		
	\subsubsection{DC 19} % *DC19
	
		\begin{compactdesc}
		
		\item[31$\vert$]
		\begin{compactdesc}
		\item[And of tenets thou shalt not talk] 1828 definition of ``tenet'': ``Any opinion, principle, dogma or doctrine which a person believes or maintains as true; as the tenets of Plato or of Cicero. The tenets of christians are adopted from the Scriptures; but different interpretations give rise to a great diversity of tenets.''
			That helps explain a bit of what is going on here.
		\end{compactdesc}
		
		\end{compactdesc}
		
	\subsubsection{DC 20} % *DC20
	
		There have been many, many, many changes to this revelation over time. In addition to that, Oliver Cowdery wrote a similar document that this one replaced. See HDDC 286--301, which also includes Cowdery's document; see JSP-D1:116--120 for the background of the revelation, and also JSP-D1:368--377 for Cowdery's document (includes a few photographs).
		See also MSDC section 20, which has a bit about the controversy over the baptismal requirements, and a couple references to the primary sources.
	
		\begin{compactdesc}
		
		\item[35$\vert$]
		\begin{compactdesc}
		\item[according to the revelations of John] JSP-D1:123 has, ``...and agreeable to the revelations of Jesus Christ which was signified by his angel unto John.'' Compare JSR 64, ``agreeable to the revelations of John.'' HDDC says that nearly all old manuscripts have ``agreeable.'' 
		\end{compactdesc}
		
		\item[37$\vert$]
		\begin{compactdesc}
		\item[Manner of baptism] As given in JSR, this verse appears just as a brief reference to Mor 6:1--4, which gives the requirements of baptism in a way similar to how they appear now in verse 37 (JSR 64). The JSR copy of the revelation is from Zebedee Coltrin's journal, and according to page 62 footnote 8, was copied in 1832. The JSR facsimile book has the text of 37 (81); so does JSP-D1:123.
		\item[Spirit of Christ] JSP-D1:123 has ``gift''; see note 61.
		\item[shall be received by baptism into his church] Makes the practice of rebaptism in the early days even more beautiful. I sort of wish that that was still done; going down into the water is a much more powerful symbol than taking the sacrament, but that just means that we have to give extra effort to prepare ourselves and renew our covenants with a firm mental commitment. Otherwise the rebaptism isn't going to help either.
		\end{compactdesc}
		
		\item[38$\vert$]
		\begin{compactdesc}
		\item[an apostle is an elder] See JSP-D1:123 note 63 for the use of this word; very different from how we understand this term in the church today. See also JSP-D1:143--144.
		\end{compactdesc}
		
		\end{compactdesc}

	\subsubsection{DC 42} % *DC42
	
		\begin{compactdesc}
		
		\item[48$\vert$]
		\begin{compactdesc}
		\item[appointed unto death] See Job 7:1, 3; 14:5, 13--14; Ps 79:11; Ps 102:19--20; % *APPOINTED
		
		\end{compactdesc}
		
		\end{compactdesc}

	\subsubsection{DC 45} % *DC45
	
		\begin{compactdesc}
		
		\item[3--5$\vert$]
		\begin{compactdesc}
		\item[Father responds because of Christ] Other examples of this:
			\begin{compactitem}
			\item 3N 27:7: ``and ye shall call upon the Father in my name that he will bless the church for my sake.''
			\item Al 33:11, 13: ``And thou didst hear me because of mine afflictions and my sincerity; and it is because of thy Son that thou hast been thus merciful unto me, therefore I will cry unto thee in all mine afflictions, for in thee is my joy; for thou hast turned thy judgments away from me, because of thy Son...ye must believe what Zenos said; for, behold he said: Thou hast turned away thy judgments because of thy Son.''
			\item Al 33:16: ``Thou art angry, O Lord, with this people, because they will not understand thy mercies which thou hast bestowed upon them because of thy Son.''
			\end{compactitem}
		\end{compactdesc}
		
		\end{compactdesc}

	\subsubsection{DC 46} % *DC46
	
		\begin{compactdesc}
		
		\item[13--14$\vert$]
		\begin{compactdesc}
		\item[to some...to others...] See also:
			\begin{compactdesc}
			\item[1N 2:16--17] (Note that it happens twice in these verses: Lehi/Nephi, then Nephi/Sam.) And it came to pass that I, Nephi, being exceedingly young, nevertheless being large in stature, and also having great desires to know of the mysteries of God, wherefore, I did cry unto the Lord; and behold he did visit me, and did soften my heart that I did believe all the words which had been spoken by my father; wherefore, I did not rebel against him like unto my brothers. 17 And I spake unto Sam, making known unto him the things which the Lord had manifested unto me by his Holy Spirit. And it came to pass that he believed in my words.
			\item[2N 31:16--17] And now, my beloved brethren, I know by this that unless a man shall endure to the end, in following the example of the Son of the living God, he cannot be saved. 17 Wherefore, do the things which I have told you I have seen that your Lord and your Redeemer should do; for, for this cause have they been shown unto me, that ye might know the gate by which ye should enter. For the gate by which ye should enter is repentance and baptism by water; and then cometh a remission of your sins by fire and by the Holy Ghost.
			\item[Jac 1:5--6] For because of faith and great anxiety, it truly had been made manifest unto us concerning our people, what things should happen unto them. 6 And we also had many revelations, and the spirit of much prophecy; wherefore, we knew of Christ and his kingdom, which should come.
			\item[Al 19:13, 31] (Lamoni sees his redeemer, then people believe on his words later.)
			\item[3N 19:20--21, 23] Father, I thank thee that thou hast given the Holy Ghost unto these whom I have chosen; and it is because of their belief in me that I have chosen them out of the world. 21 Father, I pray thee that thou wilt give the Holy Ghost unto all them that shall believe in their words...23 And now Father, I pray unto thee for them, and also for all those who shall believe on their words, that they may believe in me, that I may be in them as thou, Father, art in me, that we may be one.
			\item[3N 19:28] Father, I thank thee that thou hast purified those whom I have chosen, because of their faith, and I pray for them, and also for them who shall believe on their words, that they may be purified in me, through faith on their words, even as they are purified in me.
			\item[Mor 7:31--32] (Note the word ``residue'' in this verse, which in the 1828 means ``That which remains after a part is taken, separated, removed or designated.'' So this may be a little tension against the idea in 46:13--14 that it doesn't say ``the rest.'' Although that scripture still stands as it's written, so it's not like these verses ruin that interpretation.) And the office of their ministry is to call men unto repentance, and to fulfil and to do the work of the covenants of the Father, which he hath made unto the children of men, to prepare the way among the children of men, by declaring the word of Christ unto the chosen vessels of the Lord, that they may bear testimony of him. 32 And by so doing, the Lord God prepareth the way that the residue of men may have faith in Christ, that the Holy Ghost may have place in their hearts, according to the power thereof; and after this manner bringeth to pass the Father, the covenants which he hath made unto the children of men.
			\item[DC 5:11--14] And in addition to your testimony, the testimony of three of my servants, whom I shall call and ordain, unto whom I will show these things, and they shall go forth with my words that are given through you. 12 Yea, they shall know of a surety that these things are true, for from heaven will I declare it unto them. 13 I will give them power that they may behold and view these things as they are; 14 And to none else will I grant this power, to receive this same testimony among this generation, in this the beginning of the rising up and the coming forth of my church out of the wilderness---clear as the moon, and fair as the sun, and terrible as an army with banners.
			\end{compactdesc}
		\end{compactdesc}
		
		\end{compactdesc}
		
	\subsubsection{DC 50} % *DC50
	
		HDDC 643--648 has an excellent background to this section.
	
		\begin{compactdesc}
		
		\item[24$\vert$]
		\begin{compactdesc}
		\item[brighter and brighter until the perfect day] See Pro 4:18: ``But the path of the just is as the shining light, that shineth more and more unto the perfect day.''
		\end{compactdesc}
		
		\end{compactdesc}

	\subsubsection{DC 64} % *DC64
	
		\begin{compactdesc}
		
		\item[34$\vert$]
		\begin{compactdesc}
		\item[requireth the heart and a willing mind] See also 1Ch 28:9, 21; 1Ch 29:1--9, 17; 2Ch 29:31; 
		\end{compactdesc}
		
		\end{compactdesc}
		
	\subsubsection{DC 68} % *DC68
	
		\begin{compactdesc}
		
		\item[25$\vert$]
		\begin{compactdesc}
		\item[teach them not to understand the doctrine...] See 4N 1:38--39; Neh 8:2, 7--8, 12--13;
		\end{compactdesc}
		
		\end{compactdesc}
		
	\subsubsection{DC 76} % *DC76
	
		W. W. Phelps is most likely the author of the poetic paraphrase; see Michael Hicks, ``Joseph Smith, W. W. Phelps, and the Poetic Paraphrase of `The Vision','' in the \textit{Journal of Mormon History} 20:2 (1994), pp. 63--84.
			Though note that it was published over Smith's name! February 1 1843 issue of \textit{Times and Seasons}.
			
		Brigham Young's reaction:
			``I refer to myself for an instance; I never could be persuaded that God would send every person to a lake of fire and brimstone, to be tormented by the Devil, to all eternity, for any little sin he might commit,---which was the doctrine handed down. After all, my traditions were such, that when the Vision came first to me, it was directly contrary and opposed to my former education. I said, Wait a little. I did not reject it; but I could not understand it. I then could feel what incorrect tradition had done for me. Suppose all that I have ever heard from my priest and parents---the way they taught me to read the Bible---had been true, my understanding would be diametrically opposed to the doctrine revealed in the Vision. I used to think and pray, to read and think, until I knew and fully understood it for myself, by the visions of the Holy Spirit. At first it actually came in contact with my own feelings, though I never could believe like the mass of the Christian world around me; but I did not know how nigh I believed, as they did. I found, however, that I was so nigh, I could shake hands with them any time I wished'' (CDBY 582).
		HDDC also has a nice discussion of the problems people had in accepting the revelation.
			
		Philo Dibble's report of the occasion, published in the \textit{Juvenile Instructor}, 1892, 303--304, and quoted in HDDC PDF 976 (for more information on the Dibble account, see JSP-D2:180--182 and note 202; he gave the account at least three times that we have recorded, all three times coming at least 45 years after the vision actually occurred, and he gave a few different details in each, including saying in 1877 that it was a ``glorious vision of the creation \&c''; the quote below from HDDC is the entire quote from Juvenile Instructor):
			The vision which is recorded in the Book of Doctrine and Covenants was given at the house of "Father Johnson." in Hvrum. Ohio, and during the time that Joseph and Sidney were in the spirit and saw the heavens open, there were other men in the room, perhaps twelve, among whom I was one during a part of the time---probably two-thirds of the time---I saw the glory and felt the power, but did not see the vision.
			The events and conversation, while they were seeing what is written (and many things were seen and related that are not written,) I will relate as minutely as is necessary.
			Joseph would, at intervals, say: "What do I see?" as one might say while looking out the window and beholding what all in the room could not see. Then he would relate what he had seen or what he was looking at. Then Sidney replied, "I see the same." Presently Sidney would say "what do I see?" and would repeat what he had seen or was seeing, and Joseph would reply, "I see the same."
			This manner of conversation was repeated at short intervals to the end of the vision, and during the whole time not a word was spoken by any other person.
			Not a sound nor motion made by anyone but Joseph and Sidney, and it seemed to me that they never moved a joint or limb during the time I was there, which I think was over an hour, and to the end of the vision.
			Joseph sat firmly and calmly all the time in the midst of a magnificent glory, but Sidney sat limp and pale, apparently as limber as a rag, observing which, Joseph remarked, smilingly, "Sidney is not used to it as I am."
		Another account, from the book \textit{Scenes in Early Church History}, 1882, page 81 (in the diaries/histories folder):
			On a subsequent visit to Hiram, I arrived at Father Johnson's just as Joseph and Sidney were coming out of the vision alluded to in the Book of Doctrine and Covenants, in which mention is made of the three glories. Joseph wore black clothes, but at this time seemed to be dressed in an element of glorious white, and his face shone as if it were transparent, but I did not see the same glory attending Sidney. Joseph appeared as strong as a lion, but Sidney seemed as weak as water, and Joseph, noticing his condition smiled and said, ``Brother Sidney is not as used to it as I am.''
		The third account is from early, 1877, in the minutes of a Payson Stake meeting. I have requested the CHL to digitize this record me.
			
		JS later remarked, in 1843:
			Paul ascended into the third heavens, and he could understand the three principal rounds of Jacob's ladder---the telestial, the terrestrial, and the celestial glories or kingdoms, where Paul saw and heard things which were not lawful for him to utter. I could explain a hundred fold more than I ever have of the glories of the kingdoms manifested to me in the vision, were I permitted, and were the people prepared to receive them (TPJS 305).
	
		\begin{compactdesc}
		
		\item[5$\vert$]
		\begin{compactdesc}
		\item[delight to honor] See Est 6.
		\end{compactdesc}
		
		\item[24$\vert$] % *DC76:24
		\begin{compactdesc}
		\item[begotten sons and daugthers unto God] This has always been interesting to me---\textit{unto} God, not \textit{of} God. If we human beings are created by some natural processes, then you could say that we are begotten (by that process) sons and daughters FOR God, or that we BELONG to God because of this generation process. 
		\end{compactdesc}
		
		\item[31$\vert$] % *DC76:31
		\begin{compactdesc}
		\item[partakers of God's power] 1Co 9:12 uses a similar phrase, ``If others be partakers of this power over you.'' There the word is \textgreek{ἐξουσία}.
		\end{compactdesc}
		
		\item[33$\vert$]
		\begin{compactdesc}
		\item[they are vessels of wrath] See Job 21:20.
		\end{compactdesc}
		
		\item[40$\vert$]
		\begin{compactdesc}
		\item[this is the gospel, the glad tidings] ``Gospel'' just means ``good news,'' so ``glad tidings'' is just a re-statement of that.
		\end{compactdesc}
		
		\item[57$\vert$]
		\begin{compactdesc}
		\item[after the order of Enoch] % *ORDEROFENOCH *ORDERENOCH *ENOCHORDER
			See JST additions to Gen 14, where we read about Melchizedek. In particular this part: ``And thus, having been approved of God, he was ordained a high priest after the order of the covenant which God made with Enoch, it being after the order of the Son of God; which order came, not by man, nor the will of man; neither by father nor mother; neither by beginning of days nor end of years; but of God; and it was delivered unto men by the calling of his own voice, according to his own will, unto as many as believed on his name. For God having sworn unto Enoch and unto his seed with an oath by himself; that every one being ordained after this order and calling should have power, by faith, to break mountains, to divide the seas, to dry up waters, to turn them out of their course...''
				Note how early this is, late 1830. Very early for these ideas about orders and being in the presence of God. Very important.
				Other uses of ``order of Enoch'':
					JST Gen 9:9: ``And I, behold, I will establish my covenant with you, which I made unto your father Enoch, concerning your seed after you...''
						See also 11, 16, 13:14, and all the verses in Moses about Enoch (those will be important for getting at the order of Enoch); 
					CDBY: He concluded by saying when we are fully organized in the Order of Enoch, and partake of its spirit, we will bring everything we have on earth, and deed it over to our Father who is in Heaven.
					CDBY: We now want to organize the Latter-day Saints, every man, woman and child among them, who has a desire to be organized, into this holy order. You may call it the Order of Enoch, you may call it co-partnership, or just what you please. It is the United Order of the Kingdom of God on the earth; but we say the Order of Enoch on the same principle you find in the revelation concerning the Priesthood, which, to avoid the too frequent repetition of the name of the Deity, is called the Priesthood after the order of Melchisedec. This order is the order of heaven, the family of heaven on the earth; it is the children of our Father here upon the earth organized into one body or one family, to operate together.
		\end{compactdesc}
		
		\end{compactdesc}

	\subsubsection{DC 84} % *DC84
	
		\begin{compactdesc}
		
		\item[24$\vert$]
		\begin{compactdesc}
		\item[swore that they should not enter into his rest] See Ps 95:8--11, especially 11.
		\end{compactdesc}
		
		\item[33$\vert$]
		\begin{compactdesc}
		\item[unto the renewing] See Isa 40:31--41:1.
			See also DC 24:12, ``And at all times, and in all places, he shall open his mouth and declare my gospel as with the voice of a trump, both day and night. And I will give unto him strength such as is not known among men.''
		\end{compactdesc}
		
		\item[45$\vert$]
		\begin{compactdesc}
		\item[Original text] starting a new paragraph: ``for the word of the Lord is truth and whatsoever is truth is light, and whatsoever is light is spirit even the spirit of Jesus Christ'' (\textit{The Joseph Smith Revelations}, 214).
		\end{compactdesc}
		
		\item[48$\vert$]
		\begin{compactdesc}
		\item[the Father teacheth him] Specifically interested here in how it's the Father doing the teaching. See Morm 5:17; Isa 2:3; 
		\end{compactdesc}
		
		\end{compactdesc} 

	\subsubsection{DC 88} % *DC88
	
		On January 11, 1833, Joseph Smith wrote to W. W. Phelps about this revelation: ``I send you the `olive leaf' which we have plucked from the tree of Paradise, the Lord's message of peace to us; for though our brethren in Zion indulge in feelings towards us, which are not according to the requirements of the new covenant, yet, we have the satisfaction of knowing that the Lord approves of us, and has accepted us, and established His name in Kirtland for the salvation of the nations; for the Lord will have a place whence His word will go forth, in these last days, in purity; for if Zion will not purify herself, so as to be approved of in all things, in His sight, He will seek another people; for His work will go on until Israel is gathered, and they who will not hear His voice, must expect to feel His wrath. Let me say unto you, seek to purify yourselves, and also the inhabitants of Zion, lest the Lord's anger be kindled to fierceness'' (TPJS 18. In TPJS this letter is dated January 14, but in JSP the date is January 11).
	
		\begin{compactdesc}
		
		\item[7--10$\vert$]
		\begin{compactdesc}
		\item[sun, moon, stars, earth] Christ is the light and power of all these spheres; he has some intimate relationship with them. See Hel 14:20--21---when Christ dies, these spheres ``refuse to give'' their light, and the earth shakes and trembles. Or could it be that with Christ dead, the power that runs them is weakened, and so they aren't able to provide light? And maybe with Christ coming into the physical world, there was more light than normal (3N 1:15, 19)? See 88:87 as well, maybe the ``refuse to give'' part shouldn't be literal. But see also 3N 8:21--22, but also the entire chapter 8, showing the connection between Christ and his creations.
		\end{compactdesc}
		
		\item[13$\vert$]
		\begin{compactdesc}
		\item[the light...which giveth life to all things] See Eth 3:14.
		\end{compactdesc}
		
		\item[15$\vert$]
		\begin{compactdesc}
		\item[spirit and the body are the soul of man] This revelation was given in December 1832; for earlier occurrences of this idea, see 2N 9:13, Al 40:11, 
		\end{compactdesc}
		
		\item[32$\vert$]
		\begin{compactdesc}
		\item[to enjoy that which they are willing to receive] We'll get whatever reward we feel most comfortable receiving; see Morm 9:1--5, especially 4.
		\end{compactdesc}
		
		\item[34$\vert$]
		\begin{compactdesc}
		\item[governed by law...] See also Hel 5:3, ``they could not be governed by the law nor justice, save it were to their destruction.''
		\end{compactdesc}
		
		\item[42$\vert$]
		\begin{compactdesc}
		\item[given a law to all things] See Ps 8:5, ``When I consider thy heavens, the work of thy fingers, the moon and the stars, which thou hast ordained...''
		\end{compactdesc}
		
		\item[66$\vert$]
		\begin{compactdesc}
		\item[Original text] behold that which you hear is as the voice of one crying in the wilderness, In the wilderness, because you cannot see him; my voice, becuase my voice is spirit, my spirit is truth, truth abideth and hath no end, and if it be in you it shall abound, (JSR 225)
		\end{compactdesc}
		
		\end{compactdesc}
		
	\subsubsection{DC 93} % *DC93
	
		\begin{compactdesc}
		
		\item[19$\vert$]
		\begin{compactdesc}
		\item[know how to worship]%
			\begin{compactdesc}
			\item[1N 15:14] And at that day shall the remnant of our seed know that they are of the house of Israel, and that they are the covenant people of the Lord; and then shall they know and come to the knowledge of their forefathers, and also to the knowledge of the gospel of their Redeemer, which was ministered unto their fathers by him; wherefore, they shall come to the knowledge of their Redeemer and the very points of his doctrine, that they may know how to come unto him and be saved.
			\item[TPJS 11--12] Search the scriptures---search the revelations which we publish, and ask your Heavenly Father, in the name of His Son Jesus Christ, to manifest the truth unto you, and if you do it with an eye single to His glory nothing doubting, He will answer you by the power of His Holy Spirit. You will then know for yourselves and not for another. You will not then be dependent on man for the knowledge of God; nor will there be any room for speculation. No; \textbf{for when men receive their instruction from Him that made them, they know how He will save them}.
			\end{compactdesc}
		\end{compactdesc}
		
		\item[30$\vert$]
		\begin{compactdesc}
		\item[independent in that sphere]%
			\begin{compactdesc}
			\item[Orson Pratt (?), Millennial Star 18:146--147, 1856] Among His numerous spiritual children, God saw that there was a great variety of capacities, that some were more noble than others; the more noble ones He set apart to be rulers in His kingdom at their appointed times. These have had power to bestow the Priesthood, and administer salvation to those of their brethren who would be obedient to their authority. They have appeared on the earth at the times before appointed them in the councils of heaven, \textbf{and in their sphere their power is absolute and independent}...\textbf{They are as independent in the sphere assigned them, as God is in His}.
			\end{compactdesc}
		\end{compactdesc}
		
		\item[36$\vert$]
		\begin{compactdesc}
		\item[glory of God is intelligence] Some quotes:
			\begin{compactdesc}
			\item[CDBY 76] But before I undertake to explain or give correct views upon this important subject, I would say to all those who are satisfied with all the knowledge they have, and want no more: to you I do not expect to be an apostle this day; but for those who are hungering and thirsting after righteousness, I pray, that they may be filled and satisfied with \textbf{the intelligence of God, even his glory}.
			\end{compactdesc}
		\end{compactdesc}
		
		\end{compactdesc}
		
	\subsubsection{DC 121} % *DC121
	
		\begin{compactdesc}
		
		\item[43$\vert$]
		\begin{compactdesc}
		\item[with sharpness] Alan Erdmann once said in a priesthood meeting that he thought this means quickly, at the exact time, or about the exact matter to be dealt with and nothing else. I like that but think it probably means more too. See Mor 9:4, ``when I speak the word of God with sharpness they tremble and anger against me; and when I use no sharpness they harden their hearts against it.''
		\end{compactdesc}
		
		\end{compactdesc}
		
	\subsubsection{DC 122} % *DC122
	
		I haven't checked over this whole section yet, but verse 9 has some clear parallels to the language in Job 14.
		
		\begin{compactdesc}
	
		\item[9$\vert$]
		\begin{compactdesc}
		\item[their bounds are set, they cannot pass] See Job 14:5.
		\item[days known/years not numbered less] See Job 14:5, 13--14.
		\end{compactdesc}
		
		\end{compactdesc}
		
	\subsubsection{DC 128} % *DC128
	
		\begin{compactdesc}
		
		\item[20$\vert$] % *DC128:20
		\begin{compactdesc}
		\item[the voice of Michael...] See this email I wrote to my mom on 1/11/2017:
			\begin{quote}
				On May 1 1842 JS's journal says,

				"I preached in the grove on the keys of the Kingdom...The keys are certain signs and words by which false spirits and personages may be detected from true, which cannot be revealed to the Elders till the Temple is completed -The rich can only get them in the Temple -the poor may get them on the Mountain top as did Moses..."

				Perhaps JS is referring to the actual book of Moses here, where "Moses was caught up into an exceedingly high mountain." Satan appears in Moses 1 and Moses learns how to tell the difference between Satan and God: "I can judge between thee and God"; "his glory has been upon me, wherefore I can judge between him and thee."

				JS translates/receives Moses 1 in June 1830, and leaves for Ohio near the beginning of 1831. So first of all it's possible that he translates Moses 1 just before, very near to, or just after having this experience with Michael and Satan on the river. JS learns the same thing that Moses did, which is to discern spirits.

				Second, the experience with Michael seems to be one of the occasions in which JS receives keys and knowledge in preparation for the endowment. BY later says, and we hear something like this in the endowment today, "Your endowment is, to receive all those ordinances in the House of the Lord, which are necessary for you, after you have departed this life, to enable you to walk back to the presence of the Father, passing the angels who stand as sentinels, being enabled to give them the key words, the signs and tokens, pertaining to the Holy Priesthood, and gain your eternal exaltation in spite of earth and hell."

				So we need the knowledge of detecting spirits, apparently, in order to get to the celestial kingdom. Michael/Adam appears to JS sometime in 1830 or so in order to give a bit of this knowledge. JS expanded on some of this a few years before giving the endowment, on June 27 1839:

				"In order to detect the devel when he transforms himself nigh unto an angel of light. When an angel of God appears unto man face to face in personage \& reaches out his hand unto the man \& he takes hold of the angels hand \& feels a substance the Same as one man would in shaking hands with another he may then know that it is an angel of God, \& he should place all Confidence in him Such personages or angels are Saints with there resurrected Bodies, but if a personage appears unto man \& offers him his hand \& the man takes hold of it \& he feels nothing or does not sens any substance he may know it is the devel, for when a Saint whose body is not resurrected appears unto man in the flesh he will not offer him his hand for this is against the law given him \& in keeping in mind these things we may detec the devil that he decieved us not."

				This is similar to DC 129, though that section comes after the endowment in 1843. But sometime in the second half of 1839 JS also teaches:

				"The Priesthood is an everlasting principle \& Existed with God from Eternity \& will to Eternity, without beginning of days or end of years. the Keys have to be brought from heaven whenever the Gospel is sent.---When they are revealed from Heaven it is by Adams Authority."

				Since Adam is Michael, it could be that the experience on the river with the devil was a step in Adam's revealing the keys of the priesthood, in this case doing it himself rather than delegating it through his authority.
			\end{quote}
		\end{compactdesc}
		
		\end{compactdesc}
		
	\subsubsection{DC 130} % *DC130
	
		\begin{compactdesc}
		
		\item[9$\vert$]
		\begin{compactdesc}
		\item[will be made like unto crystal] See Job 37:18, ``Hast thou with him spread out the sky, which is strong, and as a molten looking glass?''
		\end{compactdesc}
		
		\end{compactdesc}
		
	\subsubsection{DC 131} % *DC131
	
		The source for this section is William Clayton's journal; see \textit{Intimate Chronicle} 102--104.
	
		\begin{compactdesc}
		
		\item[4$\vert$] % *DC131:4
		\begin{compactdesc}
		\item[``increase'' as children] How do we know that ``increase'' in this verse is referring to children? We know from Clayton's journal, which reports the following before the comments in section 131:
			``He said that except a man and his wife enter into an everlasting covenant and be married for eternity while in this probation by the power and authority of the Holy priesthood they will cease to increase when they die (i.e. they will not have any children in the resurrection), but those who are married by the power and authority of the priesthood in this life and continue without committing the sin against the Holy Ghost will continue to increase and have children in the celestial glory'' (102).
		\end{compactdesc}
		
		\end{compactdesc}
		
	\subsubsection{DC 132} % *DC132
	
		See ``polygamy.''
	
		For the date and other background information, see HC 5:xxix--xxxiv; HDDC 1731--1737; Orson Pratt, JD 13:192--193; Intimate Chronicle 110; 
		
		JS dictated the revelation to William Clayton. Note that in Clayton's 1874 statement about the dictation, he says, ``Joseph and Hyrum then sat down and Joseph commenced to dictate the revelation on celestial marriage, and I wrote it, sentence by sentence, as he dictated. After the whole was written, Joseph asked me to read it through, slowly and carefully, which I did, and he pronounced it correct. \textbf{He then remarked that there was much more that he could write on the same subject, but what was written was sufficient for the present}.''
		
		Joseph F. Smith, in JD 20:29, said: ``When the revelation was written, in 1843, it was for a special purpose, by the request of the Patriarch Hyrum Smith, and was not then designed to go forth to the church or to the world. It is most probable that had it been then written with a view to its going out as a doctrine of the church, it would have been presented in a somewhat different form. There are personalities contained in a part of it which are not relevant to the principle itself, but rather to the circumstances which necessitated its being written at that time. Joseph Smith, on the day it was written, expressly declared that there was a great deal more connected with the doctrine which would be revealed in due time, but this was sufficient for the occasion, and was made to suffice for the time. And, indeed, I think it much more than many are prepared to live up to even now.''
		
		\begin{compactdesc}
		
		\item[17$\vert$] % *DC132:17
		\begin{compactdesc}
		\item[separately and singly] JS uses this phrase in speaking to William Clayton about children on May 18, 1843: ``I asked the President w[h]ether children who die in infancy will grow. He answered ``no, we shall receive them precisely in the same state as they died i.e. no larger. They will have as much intelligence as we shall but shall always remain separate and single. They will have no increase. Children who are born dead will have full grown bodies being made up by the resurrection'' (Clayton's journal, 104).
		\end{compactdesc}
		
		\item[26$\vert$] % *DC132:26
		\begin{compactdesc}
		\item[a free pass for sin?] Below is the text of an email I wrote to Tyler Ransom on December 7, 2016, when he asked about this verse.
				%%%%%%%%%%%%%%%%%%%%%%%%%
					I've put a quote below from Joseph Fielding Smith that's a decently direct answer to the question of what is going on in this verse. His answer is that we have to supply the required repentance, so that no one gets a free pass--you might have been sealed but unless you repent of your sins you aren't going to be able to enter into the kingdom. In other words the material here in verse 26 is not saying that there's a repentance-free way to get into the kingdom; everyone has to repent, no matter what. JFS also calls attention to how awful the punishments here described really will be.

					I did find one person who suggested something interesting: "It seems as if the Lord spoke verse 26 specifically to set Joseph's mind at ease about Emma's eternal destiny." She had understandable trouble with the revelation and its commandments, and reading 26 as a remark about Emma is definitely a gentler reading than the fire-and-brimstone one that JFS gives below. Emma vacillated and went back and forth but eventually on September 28 1843, not long after 132 was written down, they are anointed with the second anointing to be king and priest, queen and priestess. 

					Later, and I didn't know this before, Emma wrote a blessing for herself at Joseph's request with the understanding that he would "sign" it or seal the blessings she had asked for. This was right before he was killed apparently. I've attached a copy of it here in case you are interested, it's quite striking. In accordance with what we were saying earlier in this thread, Emma asks for "the Spirit of God to know and understand myself, that I might be able to overcome whatever of tradition or nature that would not tend to my exaltation in the eternal worlds." Later she says, "I desire with all my heart to honor and respect my husband as my head, ever to live in his confidence and by acting in unison with him retain the place which God has given me by his side." To me these are indications that she had humble and sincere desires to do what was required of her, even if she sometimes didn't do it and even if it was exactly the thing she didn't want to do; this seems like a reasonable example of the attitude called for in verse 26, the only real attitude that the sealing will allow.

					Some questions in my mind still remain, though, such as about the sealing by the HSOP--when and by what means that happens. But this is part of the whole Nauvoo priesthood stuff that is extraordinarily interesting but also just a complete maze when you try to start getting your head around it.

					Joseph Fielding Smith, Doctrines of Salvation, volume 2:

					"Verse 26, in section 132, is the most abused passage in any scripture. The Lord has never promised any soul that he may be taken into exaltation without the spirit of repentance. While repentance is not stated in this passage, yet it is, and must be, implied. It is strange to me that everyone knows about verse 26, but it seems that they have never read or heard of Matthew 12:31-32, where the Lord tells us the same thing in substance as we find in verse 26, section 132.

					It is wrong to take one passage of scripture and isolate it from all other teachings dealing with the same subject. We should bring together all that has been said by authority on the question. If we were to make a photograph, it would be necessary for all of your rays of light to be focused properly on the subject. If this were not done then a blurred picture would be the result. This is the case when we try to obtain a mental picture, when we have only a portion of the facts dealing with the subject we are considering. Therefore we must find out what else has been said about salvation.

					The Lord said by his own mouth: "And he that endureth not unto the end, the same is he that is also hewn down and cast into the fire, from whence they can no more return, because of the justice of the Father. And this is the word which he hath given unto the children of men. And for this cause he fulfilleth the words which he hath given, and he lieth not, but fulfilleth all his words. And no unclean thing can enter into his kingdom; therefore nothing entereth into his rest save it be those who have washed their garments in my blood, because of their faith, and the repentance of all their sins, and their faithfulness unto the end.

					So we must conclude that those spoken of in verse 26 are those who, having sinned, have fully repented and are willing to pay the price of their sinning, else the blessings of exaltation will not follow. Repentance is absolutely necessary for the forgiveness, and the person having sinned must be cleansed.

					John said: "There is a sin unto death." "If any man see his brother sin a sin which is not unto death, he shall ask, and he shall give him life for them that sin not unto death. There is a sin unto death: I do not say that he shall pray for it."

					The Lord, in verse 27, has pointed out some sins unto death for which there is no forgiveness. It will do no good for one to pray for his brother for forgiveness from such a sin. All other sins, including blasphemy against the Son of God, may be forgiven men, on their true repentance. If they do not repent, then no matter what the sin may be, or the covenant violated, the guilty party or parties will never enter into the kingdom of God!

					Here is something which those who contend that the Lord has granted immunity from their sins to some, if they have received certain sealings by the Holy Spirit of promise, have overlooked in this passage. I call attention to these two things. If covenants are broken and enormous sins are committed, but not unto death, there are certain punishments to be inflicted. The mere confession is not enough; the sinners are: 1-to "be destroyed in the flesh"; and 2-to "be delivered unto the buffetings of Satan unto the day of redemption."

					Who in the world is so foolish as to wish to sin with the hope of forgiveness, if such a penalty is to be inflicted? No one but a fool! To be "destroyed in the flesh" means exactly that. We cannot destroy men in the flesh, because we do not control the lives of men and do not have power to pass sentences upon them which involve capital punishment. In the days when there was a theocracy on the earth, then this decree was enforced. What the Lord will do in lieu of this, because we cannot destroy in the flesh, I am unable to say, but it will have to be made up in some other way.

					Then to be turned over to the buffetings of Satan unto the day of redemption, which is the resurrection, must be something horrible in its nature. Who wishes to endure such torment? No one but a fool! I have seen their anguish. I have heard their pleadings for relief and their pitiful cries that they cannot endure the torment. This was in this life. Add to that, the torment in the spirit world before the redemption comes-all of this, mark you, coming after severe and humble repentance!"
				%%%%%%%%%%%%%%%%%%%%%%%%%
		\end{compactdesc}
		
		\end{compactdesc} 

\section{Pearl of Great Price}

\subsection{Moses} % *MOSES

	\subsubsection{Moses 1} % *MOS1
	
		\begin{compactdesc}
		
		\item[39$\vert$]
		\begin{compactdesc}
		\item[my work and my glory] See Ps 90:16. Ps 90 is subtitled ``a prayer of Moses the man of God.'' That entire chapter is very interesting but especially this verse, I didn't realize that it was here.
		\end{compactdesc}
		
		\end{compactdesc}

	\subsubsection{Moses 7} % *MOS7

		\begin{description}

		\item[53$\vert$]
		\begin{compactdesc}
		\item[Messiah] Note that the use of \textit{Messiah} here without an article (definite or indefinite) parallels the use in DC 13:1.
		\end{compactdesc}

		\end{description}
		
\subsection{Abraham} % *AB *ABRAHAM

	\subsubsection{Facsimile 1} % *ABFAC1
	
		\begin{compactdesc}
		
		\item[12$\vert$]
		\begin{compactdesc}
		\item[Raukeeyang] I think he's referring to the Hebrew word \he{rAqiy`a}, ``expanse.''
		\end{compactdesc}
		
		\end{compactdesc}
	
	\subsubsection{Facsimile 2} % *ABFAC2
	
		This diagram is actually a hypocephalus. From Wikipedia: ``The Joseph Smith Hypocephalus (also known as the Hypocephalus of Sheshonq) was a papyrus fragment, part of the original Joseph Smith Papyri, found in the Gurneh area of Thebes, Egypt, around the year 1818. The owner's name, Sheshonq, is found in the hieroglyphic text on said hypocephalus. Three hypocephali in the British Museum (37909, 8445c, and 8445f) are similar to the Joseph Smith Hypocephalus both in layout and text and were also found in Thebes.''
	
	\subsubsection{Facsimile 3} % *ABFAC3

	\subsubsection{Abraham 1} % *AB1
	
		\begin{compactdesc}
		
		\item[2$\vert$]
		\begin{compactdesc}
		\item[I sought for the blessings of the fathers...] Remember Cowdery's introduction to the Patriarchal Blessings book, written in September 1835, which was during the time they were working on the translation of Abraham: ``But before baptism, our souls were drawn out in mighty prayer to know how we might obtain the blessings of baptism and of the Holy Spirit, according to the order of God, and we diligently sought for the right of the fathers and the authority of the holy priesthood, and the power to admin[ister] in the same: for we desired to be followers of righteousness and the possessors of greater knowledge, even the knowledge of the mysteries of the kingdom of God...After this we received the high and holy priesthood: but an account of this will be given elsewhere, or in another place'' (EPB 3).
		\item[Relation to Egyptian papers] Note that this verse follows the order of the degrees given for Abraham's name in the Egyptian papers; see JSEP 33--34. It follows the degrees starting with the first and going to the fifth.
		\end{compactdesc}
		
		\end{compactdesc}

	\subsubsection{Abraham 3} % *AB3
	
		\begin{compactdesc}
		
		\item[22$\vert$]
		\begin{compactdesc}
		\item[the intelligences that were organized before the world was] I've never been clear on what ``were organized'' means here; most people seem to think it means that individual intelligences were at some point \textit{dis}organized, or hadn't yet taken such a form as to be called ``intelligences,'' or something like that. I imagine God gathering together the various parts and putting them in the right order so that they could be ``organized'' and calld intelligences. 
		
		But another way of reading this verse is just to say that the intelligences were organized as in, they were gathered together into various ranks and hierarchies, already taking their final form (or in their original form, which was never a disorganized form.) I need to look at how JS uses ``organize'' in contexts like this one though.
		
		See also CDBY 61, ``But Christ is the head of all for he is our Elder Brother for we were once organized. Before God and Jesus was the first born, or begotten of the Father, and we were sent here upon this earth to choose bodies and dwell in the flesh as our Father who is in heaven.'' I read this quote with different punctuation: ``But Christ is the head of all for he is our Elder Brother; for we were once organized before [i.e. in the presence of] God, and Jesus was the first born, or begotten of the Father [which is why he is the head of all], and we were sent here upon this earth to choose bodies and dwell in the flesh as our Father who is in heaven.'' In fact, this is the correct way to read it; see Eugene England's BYU Studies article ``George Laub's Nauvoo Journal'' (1978), where it gives this quote as, ``But Christ is the head of all, for he is our head \& Elder Brother. For we was once organised before God, \& Jesus was the first born or begotten of the Father \& we ware sent here upon this Earth to choose bodys \& dwell here in the flesh as our Father who is in heaven.''
		
		Beginning to lose faith in CDBY, or least I'm learning to be more careful. Very disappointing if that's really what the quote is. Tomorrow I will look more into this.
		
		At the least it seems that that quote from CDBY 61 could be reporting about the same event as in Ab 3:22. 
		\end{compactdesc}
		
		\end{compactdesc}
		
	\subsubsection{Abraham 4} % *AB4
	
		\begin{compactdesc}
		
		\item[2$\vert$]
		\begin{compactdesc}
		\item[empty and desolate] Note that in JS's new translation of Genesis in 1830 he did not change this phrase from the KJV, he kept it as ``without form, and void'' (Mos 2:2). But here in Abraham he changes it to ``empty and desolate.'' As I've just learned with the help of Marquardt's essay in JSEP, Seixas's 1834 Hebrew grammar gives ``empty'' and ``desolate'' for the \textit{tohu wbohu} Hebrew words that are in the verse in Genesis. So JS gets this translation from the grammar.
			And as Marquardt points out, one of the biggest differences between JS's translation of the book of Abraham in 1835 (1:1--2:18) and the rest (in 1842) is that he learned some things about Hebrew in 1836. That's why you get all the transliterated Hebrew words in chapter 3 and not earlier---he didn't know about it when he did the earlier stuff.
		\end{compactdesc}
		
		\end{compactdesc}

\section{Old Testament}

\subsection{OT interesting stories to teach} % *OTINTERESTING

\begin{itemize}

\item \textbf{The congregation hears the voice of the Lord (Ex 3)}

	See also Deu 4:9--15, 33, 36; 5:2--5, 22--29;

\item \textbf{Aaron's sons washed, anointed, clothed, and then the glory of the Lord is manifest (Lev 8--9; Num 8, 18;)}

	See also Num 3:5--28, 40--51; the sons of Levi are like an offering to the Lord instead of the firstborn sons, as in the passover. They are offered as a sacrifice for the children of Israel (see Num 8:11 especially). They are separated from others and offered as a sacrifice.

\item \textbf{Aaron's sons Nadab and Abihu killed for improper sacrifice (Lev 10)}

	The context is important: this happens right after the washing, anointing, clothing, and consecrating of the priests (Aaron's sons). Aaron has four sons. Two are killed, and two continue to officiate (the two left are in verse 12). Aaron laments in verse 19. They make a mistake in officiating in verse 16--18, but Moses is apparently fine with it after Aaron says that he is upset. I haven't thought about this story enough yet.

	See also Lev 16:1--2; Num 3:2--4; Num 26:60--61;

\item \textbf{Instructions for offering sacrifice for the people on the day of atonement (Lev 16)}

	See 29--31 for an introduction.

\item \textbf{Earliest warnings against pornography (Lev 18)}

	Nice scriptures here: ``Ye shall do my judgments, and keep mine ordinances, to walk therein: I am the LORD your God. Ye shall therefore keep my statutes, and my judgments: which if a man do, he shall live in them: I am the LORD'' (4--5).

\item \textbf{Son of Israelitish woman and Egyptian fights in camp, blasphemes, is stoned (Lev 24:10--16, 23)}

\item \textbf{Promises for keeping the statutes and commandments of the Lord (Lev 26:3--12)}

	The most important is in 12: ``And I will walk amon gyou, and will be your God, and ye shall be my people.'' Cf. Mos 7:69, ``And Enouch and all his people walked with God, and he dwelt in the midst of Zion.''

\item \textbf{Instructions on camp set-up for the children of Israel, in the east, south, west, and north, around the tabernacle (Num 2)}

\item \textbf{Sons of Levi given to Aaron and his sons for ministering in the priest's office (Num 3:5--28, 40--51)}

\item \textbf{Organization of the sons of Levi around the tabernacle (Num 3:29--38}

\item \textbf{Instructions to cover and transport the implements of the tabernacle (Num 4:5--15)}

\item \textbf{Instructions on not touching the holy items of the tabernacle (Num 4:15;)}

\item \textbf{Using holy water and dust in discovering and punishing adultery (Num 5:11--31)}

\item \textbf{Instructions for Nazarite vows (Num 6:1--21)}

	See Jud 14--16; also 1Sam 1;

\item \textbf{Song of the sons of Aaron (Num 6:23--27)}

\item \textbf{Dedication of the tabernacle to the Lord (Num 7)}

\item \textbf{God speaks to Moses from the throne of the mercy seat (Num 7:89)}

\item \textbf{Cloud by day, appearance of fire by night; cloud moves, people of Israel move (Num 9:15--23)}

\item \textbf{People complain, the Lord's fire consumes them, Moses quenches the fire (Num 11:1--3)}

\item \textbf{People complain about manna; Moses says it's too heavy to carry; calling of seventy; Lord promises meat; seventy prophesy; meat given, those who eat it die (Num 11:4--35)}

	Includes this verse, verse 29: ``And Moses said unto him, Enviest thou for my sake? would God that all the LORD's people were prophets, and that the LORD would put his spirit upon them!''

\item \textbf{Miriam and Aaron complain about Moses' prophethood; the Lord responds, makes Miriam leprous, and she stays away seven days (Num 12)}

\item \textbf{The people complain and wish they had stayed in Israel (Num 14:2--4)}

\item \textbf{Children of Israel complain against Moses and Aaron; they are sentenced to 40 years in the wilderness (Num 13:17--14)}

\item \textbf{A man stoned for gathering sticks on the Sabbath (Num 15:32--36)}

\item \textbf{Korah, Dathan, and Abiram killed for rebellion and trying to assume the priesthood (Num 16)}

	There are many aspects to this story, so you have to read the whole chapter. Note what tribe they come from, though---one from Levi, and two from Reuben.

	Note Num 26:9--11; Deu 11:6;

\item \textbf{Aaron's rod (of the family of Levi) blossoms, while other rods don't (Num 17)}

\item \textbf{Instructions to sacrifice the red heifer (Num 19:1--10)}

\item \textbf{Miriam dies; no water, water from the rock, Moses and Aaron cursed for disobedience (Num 20:1--13)}

	See also Num 27:12--14;

\item \textbf{Aaron dies, and Eleazar his son receives his garments on Mount Hor (Num 20:22--29)}

\item \textbf{Fiery serpents, the brass serpent, and look and live (Num 21:4--9)}

\item \textbf{Balaam instructed by the Lord; the donkey sees the angel and talks (Num 22)}

\item \textbf{Balaam talks with the Lord for Balak; he is obedient to the Lord's commands, and sees a vision of Israel's future and of Christ (Num 23--24)}

	Note verses 24:2--4 and 24:15--16, where it describes the process of receiving the vision; very familiar language in those verses.

\item \textbf{Phinehas slays the man with the Midianitish woman, and is promised an everlasting priesthood (Num 25:1--15)} % *PHINEHAS

	For more about Phinehas, see Num 31:6; Jos 22; 24:33; Jud 20:28; 1Sam 1:3, 2:34, 4:4, 11; 14:3 (same Phineas there? I think it's different); 1Ch 9:20; Ps 106:30--31;

	Phinehas is Aaron's grandson and Eleazar's son. The verses about the priesthood are 10--13.

\item \textbf{The daughters of Zelophehad ask for and receive an inheritance (Num 27:1--11)}

	See also Num 26:28--33; 36:1--12; Jos 17:3--6;

\item \textbf{Joshua set over the people; receives authority by laying on of hands (Num 27:15--23)}

\item \textbf{Posterity of Reuben and Gad want a different land to inherit; they promise to fight for it (Num 32:1--33)}

\item \textbf{Death of Aaron on Mount Hor (Num 33:37--39)}

\item \textbf{Moses recounts many things which happened in the wilderness and how the children of Israel screwed up (Deu 1--3}

\item \textbf{Israel kills everyone in the towns they conquer, including women and children (Deu 2:34, 3:6)}

\item \textbf{Moses asks to see the land beyond Jordan, but his prayer is turned away (Deu 3:23--29)}

	See also Deu 32:48--52, especially 51; Deu 33--34; Jos 5:6; 14;6--15. This could be a very powerful story: Moses was the most righteous and the leader, always trying to get the people to do good, but he was denied entry into the promised land, while everyone else from the congregation was let in (though they stayed in the wilderness 40 years so that most could die; need to look at that too). It's a very powerful story.

\item \textbf{Moses teaches doctrine, history, and the nature of God (Deu 4:1--40; 5; 6)}

	It's hard to break this down further; there is a lot of good stuff in these chapters and it's worth reading them a few times. Lots of interesting points.

\item \textbf{Talk of the word when you sit, walk, lie down, and rise up (Deu 6:4--9)}

	See also Deu 11:18--20; 17:18--19;

\item \textbf{Moses recounts his experiences on the mountain with the tables of stone, and the sins of the people (Deu 9--10)}

\item \textbf{A warning about false prophets (Deu 13:1--5)}

\item \textbf{How to treat and give to the poor (Deu 15:7--11)}

\item \textbf{Moses makes a prophecy of another prophet to be raised up among the people (Deu 18:15--19)}

	Maybe include verses 20--22 here as well. How do we know this prophet refers to Christ? I'm not sure.

\item \textbf{The Lord walks among the camp (Deu 23:14)}

\item \textbf{We are blessed in all things for keeping the word of the Lord (Deu 28:1--14)}

\item \textbf{When we don't obey, lots of bad things are going to happen (Deu 28:15--68)}

\item \textbf{Our life will hang in doubt; at night we'll wish it was morning, and in the morning we'll wish it was night (Deu 28:66--67)}

\item \textbf{Commandments not far (not in heaven or in the sea); good and evil set before us; choose good (Deu 30:11--20)}

\item \textbf{Moses prepares to die; the Lord appears at the tabernacle; Moses prepares the song and gathers his people together to hear it (Deu 31)}

\item \textbf{The song of Moses (Deu 32:1--43)}

\item \textbf{Moses' last days, and he is taken into heaven (Deu 32:48--34)}

\item \textbf{Moses' final blessings on the tribes of Israel (Deu 33:6--29)}

\item \textbf{Joshua assumes the role of prophet after the death of Moses (Jos 1)}

\item \textbf{The harlot Rahab receives the Israelite spies and covenants with them for her safety (Jos 2)}

	For her safety, see Jos 6:17, 22--25

\item \textbf{The Israelites pass over the Jordan on dry ground; Joshua takes the mantle of prophethood (Jos 3)}

	For Joshua, see also 4:14; 5:27;

\item \textbf{Joshua prepares twelve stones as a memorial; the children of Israel cross the Jordan on dry ground (Jos 4)}

\item \textbf{A sword-bearing angel appears to Joshua; he removes his shoes (Jos 5:13--15)}

\item \textbf{The children of Israel march around Jericho, blow trumpets and shout, and take the city (Jos 6)}

\item \textbf{Children of Israel lose to Ai in battle because one man took of the spoils of Jericho; they kill that man, and defeat the people of Ai (Jos 7--8:29)}

\item \textbf{Joshua builds an altar for sacrifice, then writes and reads the law to the people (Jos 8:30--35)}

\item \textbf{Joshua makes the sun and moon stand still (Jos 10:12--14)}

\item \textbf{Why Joshua and Caleb got to enter the promised land (Jos 14:6--15)}

\item \textbf{Inheritance of Joseph, Ephraim, and Manasseh, separate from the rest (Jos 16)}

\item \textbf{The children of Joseph get another inheritance (Jos 17:7--18)}

\item \textbf{Tabernacle set up in Shiloh; Shiloh becomes the governing city or region, it seems (Jos 18:1--11)}

	Note that this is the first occurrence of the word ``tabernacle'' since Deu 31.
	
\item \textbf{Joshua receives his inheritance (Jos 19:49--51)}

\item \textbf{Levites get forty-eight cities for inheritance (Jos 21:1--42)}

	See also 1Ch 6:54--81. I remember something about this from RSR but I haven't found it yet.

\item \textbf{Tribes of Reuben, Gad, and Mansseh go to their inheritance lands; they build an altar, not for sacrifice but as a witness for the future (Jos 22)} % *COVENANTAGREEMENT

	Very interesting story, with lots of temple implications. There's a nice parallel here with the Anti-Nephi-Lehis in the BOM, who make a covenant to not use weapons any longer (Al 24:18). Neither the altar-building nor the weapon thing is an ``official'' covenant or ordinance of salvation, but there seems to be some space for an ``unofficial'' covenant, a way of showing dedication and commitment to the Lord on your own. A symbolic gesture, but something more than that; something you want the Lord to see as more than that. Nice story.
	
	For more covenants like this, see 1Sam 20:11--23, 41--42; 23:16--18; 2Sam 21:7; 2Ki 11; 2Ki 23:1--3; 1Ch 11:3; 2Ch 15:12; 2Ch 23:1--3, 16; 2Ch 29:10; 2Ch 34:29--33; Hel 6:17--39; 3N 5:4; 3N 6:28--30; CDBY 43: ``And I would say the same for the Twelve, don't make a covenant to support them unless you intend to abide by their counsel; and if they do not counsel you as you please, don't turn round and oppose them. I want every man, before he enters into a covenant, to know what he is going to do; but we want to know if this people will support the priesthood in the name of Israel's God. If you say you will, do so''; DC 5:3; DC 5:26--27; DC 54:3--6; 

\item \textbf{Joshua's last sermon and instructions (Jos 23--24:28)}

\item \textbf{An angel appears to the children of Israel to rebuke them (Jud 2:1--5)}

\item \textbf{Judges established in the land (Jud 2:16--19)}

\item \textbf{Spirit comes on Othniel to deliver Israel from Chushanrishathaim (Jud 3:8--11)}

\item \textbf{The Lord raises up Ehud to deliver Israel, and he kills Eglon, a very fat king (Jud 3:12--30)}

\item \textbf{Deborah the prophetess judges Israel (Jud 4:4--14)}

\item \textbf{The song of Deborah and Barak (Jud 5)}

\item \textbf{The Lord calls Gideon to fight against the god Baal and the Midianite armies (Jud 6:11--40)}

	Lots of interesting aspects to this story. Note that the Lord calls a different prophet first, unnamed, to warn the people (7--10); Gideon asks for and receives at least three signs to be sure that the Lord is really with him; Gideon is worried about seeing the angel of the Lord face to face, but is told he won't die (see Jud 13:22);

\item \textbf{Gideon's army reduced to 300 by those who lap water; they win the battles; Gideon refuses the throne, and passes away (Jud 7--8)}

\item \textbf{Abimelech, son of Gideon, kills his brothers, becomes king, tells a parable of trees, fights in a war, gets hurt by a woman, and killed by his servant (Jud 9}

\item \textbf{The Lord vows that he will no longer deliver the children of Israel (Jud 10:10--14)}

\item \textbf{The Spirit comes on Jephthah; he successfully defends Israel, makes a stupid vow, and then has to sacrifice his only daughter (Jud 11:29--40}

\item \textbf{Ephraimites known (and killed) because they can't say shibboleth right (Jud 12:1--6}

\item \textbf{Angel of the Lord appears to Manoah and her husband to announce that she will bear a child; temple imagery}

	A very powerful and important story, a clear divine manifestation with temple imagery, sacrifice, and lots of other good things. For verse 22, about seeing God and dying, see Jud 6:22--23, where Gideon has the same worry. The temple imagery happens in multiple places but clearly with the name in verses 17--18.

\item \textbf{Samson kills a lion, wants a wife, eats honey, loses a riddle contest, slaughters people, burns fields, drinks water; then loses his strength when he loses his hair (Jud 14--16)}

	Violent story but interesting. Samson was a Nazarite ``from [his] mother's womb'' (16:17), and that's where he got the power from. When he broke his vow, he lost the power. Note especially 16:20: ``And he [Samson] wist not that the Lord was departed from him.'' He didn't know that the power of the Spirit had left him. Powerful story in many ways, I think. He gets strength one more time to kill everyone in the house after his eyes are put out, and he kills himself as well. There was much more to this story than I realized. Delilah wanted to trick him and didn't care about him; he had fallen in love with the wrong woman.
	
	Note that Timnath is where Samson sees the harlot; same place as Jos 19:50? The Philistine takeover really changed things I guess.

\item \textbf{Micah tries to set up a priesthood in his own house; is betrayed by his own priest and by the children of Dan; the priest's harlot and her death (Jud 17--19)}

	Very interesting story that I'd never heard before. Michah seems to want to be able to get the priesthood in his own house, even though he has idols and images to worship. Some good temple imagery here with the clothing and priesthood. That imagery is repeated in chapter 18, and the final verse of 18 mentions the temple. So this would be a good story to study in-depth to see if they were trying to set up a form of the priesthood or temple, or wanted the power but couldn't have it, or something. Reminds me of Pharoah in the book of Abraham. Note that there is no king in Israel during this time, which could be part of it.
	Note where this is happening---mount Ephraim, mentioned in Jos 19:50.

	Jud 19 records what is probably the worst story I've read so far in the OT, about a woman whose husband abandons her and she is raped all night and then cut into pieces. The ``certain Levite'' that is her husband seems to be the same Levite priest spoken of in Jud 17--18; I'm not sure exactly what is going on here.
	
\item \textbf{In response to the harlot's death, the people war against Benjamin, are beaten twice, then defeat him; they try to find wives for the new groups (Jud 20--21)}

	Phinehas makes an appearance in this story. Also, the people ask at the ``house of the Lord'' what they should do, and they aren't able to succeed until after they have fasted and offered the burnt and peace offerings to the Lord. That's a nice though and could be good for a lesson.
	
\item \textbf{The story of Ruth and Naomi (Ruth 1--4}

	A beautiful story. Check out the possible temple symbolism in 3:3, 15; resurrection symbolism in 4:1--11; the genealogy of David at the end of chapter 4.
	
\item \textbf{Hannah, mother of Samuel, asks for a son and promises to dedicate him to the Lord (1Sam 1)}

	A lot of this happens in the house of the Lord in Shiloh, where they go to offer sacrifice.
	
\item \textbf{Song of Hannah in praise of the Lord (1Sam 2:1--10)}

	Some nice temple language, especially in verse 10.
	
\item \textbf{Samuel grows up in the Lord; the Lord is displeased with the temple service of Eli, and talks about being called by the Lord (1Sam 2:12--36)}

	This is an interesting story---it looks like Eli, the priest at the temple in Shiloh, is incorrectly eating of the sacrifices that are offered by the people. Verses 27--36 are a rebuke by the Lord and show the importance of being called---and how the Lord will call someone else if the person in some position isn't getting the job done. Notice that the rebuke comes from a prophet as well, a ``man of God,'' and and begins with ``thus saith the Lord.''
	
\item \textbf{Samuel is called to be a prophet (1Sam 3)}

	This is a great story, especially when Samuel changes the words he says in response to the Lord's calling him. This is a pretty clear instance of a prophet being called by the Lord himself. This is all over the place in the OT.
	
	Great temple language here too; the Lord manifests himself in the temple.
	
\item \textbf{The Israelites lose the ark in battle; Eli dies when he hears the news, and the glory leaves Israel (1Sam 4)}

\item \textbf{The gods of the Philistines bow to the ark; Philistines die because of the ark; they return it to Israel with offerings, and the Israelites die who look inside the ark (1Sam 5--6)}

See also 1Ch 13.

\item \textbf{Ark is returned; Eleazar is sanctified to take care of it; the people obey the Lord and are delivered from the Philistines; Samuel preaches (1Sam 7)}

	Note that Samuel has an altar in his house, which is interesting.

\item \textbf{The people demand a king, and the Lord and Samuel acquiesce (1Sam 8)}

\item \textbf{Saul the son of Kish is sent by the Lord to Samuel to be anointed to lead the people in battle; he is anointed king, is spiritually reborn, and noted among the people. After winning a battle, he is actually made king (1Sam 9--11)}

	A very interesting story. One is that the chapter talks about the difference between a seer and a prophet---there isn't one, at least according to this. It uses the word ``seer'' a number of times, and this is the first time in the OT that the word occurs. Another interesting thing is that Saul is called even though he comes from a family of no reputation. A third thing is in the last verse of 9, where Samuel sends Saul's servant away so that he can speak to Saul privately, to show him ``the word of God.'' 9 and 10 are very interesting chapters about prophets and how they are called and what they do.
	
	In 11:14--15 Saul is actually made king, ``before the Lord'' in Gilgal.
	
\item \textbf{Saul offers a sacrifice to the Lord when he shouldn't, and Samuel tells him that his kingdom will be taken away and the Lord will seek another (1Sam 13:8--14)}

	Why wasn't it ok for Saul to offer the sacrifice? I think because he wasn't the priest in charge of that. Samuel was, I think.
	
\item \textbf{Saul commands the armies not to eat; his son Jonathan eats some honey; Saul can't get an answer from the Lord until he figures out that Jonathan ate the honey; he wants to kill Jonathan for his sin, but the people save him (1Sam 14:19--52)}

	Note that Saul calls for the ark right before this story, in verse 18.
	
\item \textbf{Saul wins a battle and takes the spoils when he shouldn't; he makes excuses, and Samuel tells him that to obey is better than sacrifice. Saul is removed as king from over the people of Israel because of his choices (1Sam 15:1--31)}

	This chapter includes the very important verse 22: ``And Samuel said, Hath the LORD as great delight in burnt offerings and sacrifices, as in obeying the voice of the LORD? Behold, to obey is better than sacrifice, and to hearken than the fat of rams.''

\item \textbf{The Lord through Samuel calls David and anoints him (1Sam 16)}

	Includes the great verse about how the Lord looks on the heart. Important for prophets and calling them.
	
\item \textbf{David kills Goliath (1Sam 17)}

	The story is actually kind of long, I hadn't ever read it before like this. There are a lot of interesting parts.
	
\item \textbf{Saul is jealous of David and comes to hate him; he tries to kill him, and eventually gives him his daughter Michal for a wife; David is with the Lord because he is obedient, but Saul has lost the spirit (1Sam 18)}

\item \textbf{Saul agrees to not kill David, but starts to try again; David escapes, and every messenger Saul sends to get David is touched by the spirit and begins to prophesy. Saul himself goes, and does the same (1Sam 19)}

\item \textbf{David leaves town to be safe; he makes a covenant with Jonathan (Saul's son), and Saul and Jonathan argue over David (1Sam 20)}

\item \textbf{David eats the holy bread, gets Goliath's sword, then pretends to be crazy to stay safe (1Sam 21)}

\item \textbf{Saul slaughters many priests in anger over David's escape; David is warned by one priest who gets away (1Sam 22)}

\item \textbf{David finds Saul in a cave, but doesn't kill him; they reconcile (1Sam 24)}

	Some temple symbolism here maybe, with the cutting of the skirt from Saul's clothes and the talk about Saul being the ``Lord's anointed.''

\item \textbf{David comes upon Saul again, but still refuses to kill him; Saul gives a song-like soliloquy, and says he won't harm David any longer (1Sam 26)}	

\item \textbf{Saul has a vision of Samuel through a medium; Samuel tells him that he has lost the kingdom through disobedience (1Sam 28)}

	Mentions the Urim. Note the footnote which says that this can't be a ``bona fide'' vision from God, because of both the medium and the effect. I'd never even heard of this story before.

\item \textbf{David's town is burned, and his people are kidnapped; he recovers them, and the spoil is divided equally back to those who lost it (1Sam 30)}

	Two things to note: David asks for the ephod to inquire something of the Lord (7--8); David requires the people to divide the spoils equally, even when the people didn't want to (22--25).
	
\item \textbf{Saul and his sons are killed (1Sam 31)}

\item \textbf{David finds out that Saul and Jonathan are dead; he mourns, then kills the man who killed Saul (1Sam 1:1--16)}

\item \textbf{David's song for Saul and Jonathan (2Sam 1:17--27)}

\item \textbf{David and Bathsheba---Uriah dies; Solomon is born; David is promised punishment (2Sam 11--12)}

\item \textbf{Micaiah, a prophet, prophecies the death of Ahab the king (1Ki 22:1--40)}

	The story is interesing on its own, but part of Micaiah's vision is seeing the Lord on his throne surrounded by the host of heaven (19--23). I'd never heard of this vision before and it seems a rarity in the OT.
	
\item \textbf{Elisha opens the eyes of one to see chariots, and the eyes of others to see Samaria (2Ki 6:8--23)} 

\item \textbf{In a famine, women fight over killing their own children to eat them (2Ki 6:24--33)} 

\item \textbf{In a famine, lepers discover an abandoned camp; the people take the spoils of the camp (2Ki 7)} 

\item \textbf{One of the children of the prophets anoints Jehu king (2Ki 9:1--13)} 

\item \textbf{Jehoiada the priest leads the people to serve God again, and anoints a new king; they kill Athaliah for treason (2Ki 11)} 

\item \textbf{Elisha prophesies a victory; he does a weird thing with arrows; he dies, and someone is brought back to life when that person touches Elisha's bones (2Ki 13:10--21)}

\item \textbf{Hezekiah prays in the temple, and Isaiah the prophet responds for the Lord; an angel of the Lord kills many people overnight (2Ki 19:14--37)}

\item \textbf{Hezekiah is appointed unto death, but prays and gets his life extended (2Ki 1--11)}

\item \textbf{During Josiah's reign, they clean the temple and find the book of the law, and they speak to Huldah the prophetess about what to do (2Ki 22)}

\item \textbf{Josiah cleanses the temple and the land in accordance with the book of the law; the Lord still tries the people because of their wickedness, and Josiah is killed, with his son taking his place (2Ki 23)}

\item \textbf{In the midst of the genealogies, we see a short prayer of Jabez to God (1Ch 4:9--10)}

\item \textbf{God helped the sons of Reuben in battle because they trusted him (1Ch 5:18--20)} 

\item \textbf{Three of David's captains break through a group of Philistines to draw water from the well at Bethlehem, and David is so moved that he can't drink the water (1Ch 11:15--19)}

\item \textbf{A couple brief verses about the unity between David and his men (1Ch 12:17--18)}

\item \textbf{They move the ark, and Uzza is slain when he tries to steady it (1Ch 13)}

\item \textbf{They bring the ark to Jerusalem and have a ceremony with songs and sacrifices; Levites are told to sanctify themselves (1Ch 15)}

\item \textbf{They continue to celebrate the ark, and David delivers a song of thanks for the occasion (1Ch 16)}

	The song definitely has some temple-like language; see verses 11, 13, 22, 27, 29.
	
\item \textbf{Nathan receives a prophecy about David, how the Lord will establish him on the throne; David responds with a psalm-like prayer (1Ch 17:3--27)}

\item \textbf{David counts the people, which he shouldn't do; the Lord punishes him, and he buys a place to make sacrifices to repent (1Ch 21)}

	Note verse 24, ``And king David said to Ornan, Nay; but I will verily buy it for the full price: for I will not take that which is thine for the LORD, nor offer burnt offerings without cost.'' It goes with 2Sam 24:24.
	
\item \textbf{David charges Solomon with the building of the temple; he makes some nice promises to him about righteousness (1Ch 22)}

\item \textbf{David gets all the people together, instructs them about the temple, distributes resources and responsibilities, and then Solomon starts to take over as the chosen of the Lord (1Ch 28--29)}

\item \textbf{The people rejoice at being able to give willingly to the temple (1Ch 29:1--9)}

\item \textbf{Solomon speaks to everyone near the tabernacle (2Ch 1:1--6)}

	Seems a bit like King Benjamin gathering the people by the temple.
	
\item \textbf{The Lord appears to Solomon and he asks for wisdom and understanding (2Ch 1:7--12)}

\item \textbf{The building and furnishing of the temple described (2Ch 2--4)}

\item \textbf{They move the ark into the temple, and the glory of the Lord fills it (2Ch 5)}

\item \textbf{Solomon's dedicatory prayer for the temple (2Ch 6)}

\item \textbf{Dedication is finished with sacrifice and the glory of the Lord appears; the Lord appears to Solomon and makes various promises (2Ch 7)} 

\item \textbf{Queen of Sheba visits Solomon; he is very wise and very wealthy (2Ch 9)}

\item \textbf{Rehoboam becomes king; the people ask for relief, and he increases their burdens (2Ch 10)}

\item \textbf{Rehoboam is wicked, and the prophet Shemaiah warns him (2Ch 12)}

\item \textbf{Asa is actually righteous, and is blessed and wins battles (2Ch 14)}

\item \textbf{Jehoshaphat goes to war with Ahab his ally; they ask the prophets to know what will happen, and many tell him he will be successful, but Michaiah dissents and reports a vision of the Lord (2Ch 18)}

\item \textbf{Excellent story about the Israelites fasting and praying at the temple for the Lord's help in battle; a prophet inspired of God promises them victory if they do certain things; they win and take many spoils (2Ch 20:1--30)}

\item \textbf{Jehoiada makes Athaliah's son king, kills Athaliah, and makes a covenant with the people to agree to be the Lord's (2Ch 23)}

\item \textbf{Joash and Jehoiada take up a collection to repair the house of the Lord; the people return to idolatry and kill prophets (2Ch 24)}

\item \textbf{Kings go to war, ignore the prophets, destroy each other, slaughter lots of people (2Ch 25)}

\item \textbf{Uzziah is righteous, but then gets prideful and wicked, and he tries to take the honor of burning incense upon himself; the people and priests resist him, and the Lord strikes him with leprosy as punishment (2Ch 26)} 

\item \textbf{War; people making improper sacrifices; they do something that the prophet says, surprisingly; but Ahaz continues in wickedness anyway with the sacrifices and everything else (2Ch 28)}

\item \textbf{Hezekiah has the priests cleanse the temple; they clean it out, offer sacrifice, invite people to the passover, and get things started again (2Ch 29--30)}

\item \textbf{They collect tithing and have an abundance (2Ch 31:5--10)}

\item \textbf{The priests recount the pride cycle of the people, as they've repeatedly turned away from the Lord, but he is merciful and forgives (Neh 9)}

\item \textbf{Esther becomes queen; she takes courage and decides to try saving her people; Haman's problem with Mordecai; (Est 1--5)}

	This is a pretty powerful and cool story, actually.
	
	Remember Haman's problem with Mordecai---he can't get over his hatred for him, and it spoils every good thing in his life.
	
\item \textbf{A long and detailed description of the children of Israel going astray after being led out of Egypt; a lot of nice verses to look at (Eze 20)}

\item \textbf{Prophets as watchmen, plus good stuff on repentance and accountability (Eze 33)}

\item \textbf{A very nice chapter about shepherds and their responsibility. (Eze 34)} 

\item \textbf{Vision of the dry bones; the sticks of Ephraim and Joseph (Eze 37)}

\item \textbf{Water flows forth from the temple to heal the world around it (Eze 47)}

\item \textbf{Daniel asks for a different diet of plainer food and water, and he ends up more healthy than the others; he also gains much more knowledge because of the way he lives (Dan 1)}

\item \textbf{The Lord shows Daniel the correct interpretation of the king's dream, which other people could not get (Dan 2)}

\item \textbf{Shadrach, Meshach, and Abed-Nego are thrown into the fire for not worshipping the golden image; they are preserved, and the Son of God walks with them (Dan 3)}

\item \textbf{Daniel interprets another dream; the king is warned (Dan 4)}

\item \textbf{Daniel's visions (Dan 7--12)}

	These are very interesting and worth reading again.
	
\item \textbf{Prophecies (Zech)}

	The book of Zechariah is actually pretty interesting. There is a lot here and it's worth reading this book again I think.
	
\end{itemize}

\subsection{JST OT} % *JSTOT

	\begin{description}
	\item[Num 16:10] Seeking not just the priesthood, but the \textit{high} priesthood.
	\item[Deu 2:30] Sihon is the one who hardens his spirit and heart, not the Lord, but does it so that might deliver him up to the children of Israel.
	\item[Deu 10:1--2] The revision makes clear that the everlasting covenant of the priesthood---must be the Melchizedek priesthood---was withheld from what Moses gave to the people. Compare to DC 84.
	\item[Deu 16:22] Don't set up a \textbf{graven} image, not just an image.
	\item[Deu 34:6] The Lord took Moses ``unto his fathers, and ``therefore no one knows of his sepulcher. The ``his'' seems to refer to Moses but it's ambiguous.
	\item[1Sam 9:16] Changed to ``to be a captain,'' from ``to be captain.'' Possibly an important office, pointing out that there could be more than one person in it? Or just making sense of what happens later, when there are other captains (I think).
	CHECKED UP TO AND INCLUDING 1SAM 11
	\end{description}
	
\subsection{2 Chronicles} % *2CH

	\subsubsection{2 Chronicles 7} % *2CH7
	
		\begin{compactdesc}
		
		\item[7$\vert$]
		\begin{compactdesc}
		\item[hallowed] The word is \he{qd+s}, ``to consecrate, sanctify, set apart.''
		\end{compactdesc}
		
		\end{compactdesc}
		
\subsection{Job} % *JOB

	\subsubsection{Job 19} % *JOB19
	
		\begin{compactdesc}
		
		\item[26$\vert$]
		\begin{compactdesc}
		\item[in my flesh shall I see God] Note the change that the BOM makes at 2N 9:4, ``in our \textit{bodies} we shall see God.''
			The Hebrew word for ``flesh'' here is H1320, \he{bA,sAr}, which appears around 269 times and is translated as ``flesh'' in over 250 of those.
		\end{compactdesc}
		
		\end{compactdesc}
		
\subsection{Isaiah} % *ISAIAH

	\subsubsection{Isaiah 2} % *ISA2
	
		\begin{compactdesc}
		
		\item[16$\vert$]
		\begin{compactdesc}
		\item[ships of Tarshish/pleasant pictures] Here are a few notes from when I was a gospel doctrine teacher.
			This is an interesting verse because 2 Nephi 12:16 has the following text:

			\begin{quote}
			And upon all the ships of the sea, and upon all the ships of Tarshish, and upon all pleasant pictures.
			\end{quote}

			But the Hebrew only has two of these phrases---the thing about Tarshish, and the thing about pictures, even though in the Hebrew it isn't about pictures (which might be good to mention too). And the LXX only has two of these phrases as well---the thing about the sea, and the last part, about pleasant ships (the Hebrew also mentions ships, not pictures). So 2N 12:16 has both of them. 

			It's clear from modern scholarship that the pleasant pictures part should be translated as ship; see some modern translations.

			Charles Thomson published an English translation of the LXX in 1808. Here's a bit from Wikipedia about that:

			\begin{quote}
			Charles Thomson's Translation of the Old Covenant is a direct translation of the Greek Septuagint version of the Old Testament into English, rare for its time. The work took 19 years to complete and was originally published in 1808. Thomson is credited with having created the work with little to no help from other scholars. Charles Thompson was a Greek scholar, and before the American Revolution, had been a teacher at several prominent schools. Thomson's translation of the entire Greek Bible, excluding the Apocrypha, was published in one-thousand sets of four volumes each, the fourth volume being Thomson's translation of the New Testament in that same year. The printer was Jane Aitken of Philadelphia.
			\end{quote}

			Brenton's translation of the LXX wasn't published until 1844, and was first published in London, so that doesn't matter here. Here's Thomson's translation of this verse, Isaiah 2:16:

			\begin{quote}
			and against every ship of the sea; and against every ensign of beauteous ships
			\end{quote}

			From a Maxwell Institute article on this:

			\begin{quote}
			Isaiah 2:16 is fully preserved on only one of these [the Dead Sea Scrolls texts found in Qumran], known as the \lq Great Isaiah Scroll.\rq \ As preserved thereon, Isaiah 2:16 is essentially the same as in the later Masoretic Text...These earlier textual witnesses thus provide no alternative information regarding the form or content of Isaiah 2:16.
			\end{quote}

			Source to go to: \textit{Journal of Book of Mormon Studies}, 14.2 (2005), 12-25, 67-71. A google search for \lq joseph smith septuagint 2 Nephi 12:16\rq \ brings it up as the first result. Pretty fair-minded discussion of the complexities. It's much more difficult to say than the footnote in the Bible and BOM make it out to be.

			However, there is no satisfactory explanation for how Joseph Smith came to include the \lq ships of the sea\rq \ part in the first place. This question is separate from that of whether we should consider this passage evidence for a divine influence on the translation. How did he know to put this in? Not clear. Every explanation is pure conjecture, requiring one to propose that he encountered a certain book at a certain time, or was introduced to a certain person who had the book at a certain time, and so on. Even then it's tenuous at best. Important to note I think, and to direct people to this fine article.
		\end{compactdesc}
		
		\end{compactdesc}

\section{New Testament} % *NT

\subsection{Matthew} % *MT

	\subsubsection{Matthew 16} % *MT16
	
		\begin{compactdesc}
		
		\item[26$\vert$]
		\begin{compactdesc}
		\item[gain the world, lose the soul] See Job 27:8.
		\end{compactdesc}
		
		\end{compactdesc}

\subsection{Luke} % *LK

	\subsubsection{Luke 2} % *LK2

		\begin{compactdesc}

		\item[52$\vert$]
		\begin{compactdesc}
		\item[favour with God and man] See also 1Sam 2:26, ``And the child Samuel grew on, and was in favour both with the LORD, and also with men.''
		\end{compactdesc}

		\end{compactdesc}
		
\subsection{John} % *JN

	\subsubsection{John 3} % *JN3
	
		\begin{compactdesc}
		
		\item[3$\vert$]
		\begin{compactdesc}
		\item[no new birth, ``cannot see the kingdom of God''] See TPJS 12, ``But except a man be born again, he cannot see the kingdom of God. This eternal truth settles the question of all men's religion. A man may be saved, after the judgment, in the terrestrial kingdom, or in the telestial kingdom, but he can never see the celestial kingdom of God, without being born of water and the Spirit.'' So either the last part of this passage is Joseph speaking loosely, or we need to add something to our interpretation of this verse in John. Or we can just view this quote from JS as talking about actually entering; he could be using ``see'' more strongly.
		\end{compactdesc}
		
		\end{compactdesc}
		
	\subsubsection{John 7} % *JN7
	
		\begin{compactdesc}
		
		\item[16$\vert$]
		\begin{compactdesc}
		\item[my doctrine is not mine] See ``doctrine of Christ.''
		\end{compactdesc}
		
		\end{compactdesc}
		
	\subsubsection{John 20} % *JN20
	
		\begin{compactdesc}
		
		\item[25$\vert$]
		\begin{compactdesc}
		\item[put my finger/thrust] See 3N 11:14.
		\end{compactdesc}
		
		\end{compactdesc}

\subsection{Romans} % *ROM

	\subsubsection{Romans 11} % *ROM11

		\begin{compactdesc}

		\item[33$\vert$]
		\begin{compactdesc}
		\item[unsearchable, past finding out] The adjectives for ``unsearchable'' and ``past finding out'' are \gr{>anexera'unhta} and \gr{>anexixn'iastoi}, respectively. The first appears a single time in the NT and the second appears twice (the other spot is Ephesians 3:8, where it is said of \gr{plo\char"007Eutos}). The second is interesting because it literally means ``not to be tracked out.'' In other words, man cannot ``track out'' or arrive at the ways of God. But this isn't an absolute statement--we can know of the ways of God through revelation; we just can't arrive at them on our own, without revelation. This is what Jacob 4:8 makes clear--it invokes the ``unsearchable'' and ``find out'' from Romans 11:33, and then makes the point without any ambiguity: ``And no man knoweth of his ways save it be revealed unto him.''

		Note that the word for ``ways'' in this verse is \gr{<odo'i}. This is the same word used in the LXX Isaiah 55:8, ``neither are your ways my ways.''
		\end{compactdesc}

		\end{compactdesc}
		
\subsection{Ephesians} % *EPH

	\subsubsection{Ephesians 2} % *EPH2
	
		\begin{compactitem}

		\item[20$\vert$]
		\begin{compactdesc}
		\item[apostles and prophets] From CDBY 55: ``The Bible says God hath set in the church, first Apostles, then comes Prophets, afterwards, because the keys and power of the Apostleship are greater than that of the Prophets.''
		\end{compactdesc}

		\end{compactitem}	

\subsection{Hebrews} % *HEB

	\subsubsection{Hebrews 3} % *HEB3

		\begin{compactitem}

		\item[6$\vert$]
		\begin{compactdesc}
		\item[firm] This is the only place that the word \textit{firm} appears in the NT of the KJV, but the Greek phrase behind it is removed in the critical edition: \gr{m'eqri t'elous beba'ian} is taken out. This is with a ranking of \textbf{B}, so fairly certain.
		\end{compactdesc}

		\end{compactitem}
		
	\subsubsection{Hebrews 5} % *HEB5
	
		\begin{compactdesc}
		
		\item[4$\vert$]
		\begin{compactdesc}
		\item[no man...but he that is called of God] See 2Ch 26:16--21, where Uzziah tries to take the honor on himself, and the Lord punishes him with leprosy.
		\end{compactdesc}
		
		\end{compactdesc}
		
\subsection{James} % *JAM

	\subsubsection{James 4} % *JAM4
	
		\begin{compactdesc}
		
		\item[17$\vert$]
		\begin{compactdesc}
		\item[Sin requires knowledge] See 3N 6:18 for the possibility of sinning in ignorance; also see the topic on sinning in ignorance under ``sin'' herein
		\item[know to do good] I'm interested in the translation of this phrase; the Greek is \gr{e>id'oti o>\cf un kal`on poie\cf in}. Specifically I'm interested in whether we should translate this construction as ``know how to do good'' or ``know to do good'' or ``know [that you should] do good'' or some other way. At some point I need to look carefully at the Greek and look at some commentaries too. See also Mt 7:11; 
		\end{compactdesc}
		
		\end{compactdesc}
		
% ----------------------
\subsection{New Testament support for JS teachings} % *NTJS *NTJOSEPHSMITH
% ----------------------

	\begin{compactdesc}
		\item[Mt 10:25] \textgreek{ἀρκετὸν τῷ μαθητῇ ἵνα γένηται ὡς ὁ διδάσκαλος αὐτοῦ, καὶ ὁ δοῦλος ὡς ὁ κύριος αὐτοῦ}.
	\end{compactdesc}